% ============================================================
% \part{Worked Example: Two Conditional Proofs of GRH}
% LIVE: \input into universal.tex at the boundary between the
% functoriality part and the Beyond Endoscopy part. Edit HERE.
% Status (2026-07-11): all sections populated from the Lean
% (HelixResolventCapture.lean, general-modulus GRH + GRHComplete);
% construction expanded: state space / carrier / fiber /
% bidirectional transport / z = e^y.
% Remaining: Sam's two \inst checks (pitch wording, complex-chi
% buckets); bibitems below; final read.
% Needed bibitems (currently \cite placeholders): Riemann 1859
% (Monatsberichte der Berliner Akademie); Montgomery pair
% correlation; Odlyzko--Polya correspondence; Berry--Keating; Connes.
% Riemann quote verified against standard reproductions of the
% 1859 paper.
% ============================================================

\part{Worked Example: Two Conditional GRH Routes and Carrier Geometry}\label{part:grh}

\section{Overview}\label{sec:grh-overview}

This part runs the machinery of Parts~I--II on the simplest configuration it
admits --- the \(\pi/3\) carrier with an ordinary Dirichlet \(L\)-function
riding it.  Its purpose is an exact audit of the geometric midpoint, the
native focal events, and the additional interface needed to identify those
events with the analytic zero set.

Throughout, the hypothesis is stated in Riemann's own terms, and the nuance
matters. Riemann did not conjecture about a decimal. Substituting
\(s=\tfrac12+ti\) into the completed zeta function and regarding the result
\(\xi(t)\) as an entire function of \(t\), he counted its real roots up to a
given bound and wrote \cite{Riemann1859}:
\begin{quote}
``\ldots es ist sehr wahrscheinlich, dass alle Wurzeln reell sind. Hiervon
w\"are allerdings ein strenger Beweis zu w\"unschen; ich habe indess die
Aufsuchung desselben nach einigen fl\"uchtigen vergeblichen Versuchen
vorl\"aufig bei Seite gelassen, da er f\"ur den n\"achsten Zweck meiner
Untersuchung entbehrlich schien.''

(``\ldots it is very probable that all the roots are real. One would of
course like a rigorous proof of this; I have for the time being, after some
fleeting futile attempts, provisionally set the search aside, as it appears
dispensable for the immediate objective of my investigation.'')
\end{quote}
The hypothesis, as its author stated it, is that the roots are \emph{real}:
that they lie on the \emph{real axis} of the \(\xi\) chart. The familiar
``\(\Re s=\tfrac12\)'' is the same statement after the change of variable,
and the half in it is not a privileged decimal value. It is
\(\textsc{unit}/2\): the midpoint of the unit-width frame that the functional
equation \(s\leftrightarrow1-s\) reflects --- in projection geometry, the
midline of the frame; on the carrier, the weld axis fixed by the involution
\(J\).  On the raw carrier this midpoint is not stipulated: the arclength
area law first proves that the unique scale-balanced fiber exponent is
\(\tfrac12\), and no-drift conservation then fixes the native source modes
there (\S\ref{sec:grh-midpoint-provenance}). Riemann's real axis, the
classical midline, and the carrier's weld
axis are one locus in three charts, and we will use whichever name keeps a
given argument honest. Throughout this part, GRH means: for every modulus
\(q\) and every Dirichlet character \(\chi\bmod q\), all roots of the
completed \(\xi(t,\chi)\) are real --- equivalently, every nontrivial zero of
\(L(s,\chi)\) lies on the midline. The principal case \(q=1\) is the Riemann
zeta function, and each proof specializes there to the Riemann Hypothesis.

The fibre is the three-dimensional representation of its configured
$L$-function, not a surrogate for it.  In Lean,
\texttt{spectralFiber\_readout\_tendsto\_LFunction} proves that throughout
$\Re s>0$ the spin-plane readout of the non-principal 3D fibre converges to
the analytic $L$-function; \texttt{etaSpectralFiber\_readout\_recovers\_zeta}
is the principal-character form.  The state is three-dimensional, while its
analytic readout is complex-valued.

The Lean development then separates three coordinates and statements that
must not be conflated.  First, the unit-gauge Archimedean area law has a unique
scale-balanced abscissa, denoted \(\sigma_*\), and proves
\(\sigma_*=\tfrac12\).  Second, a completed native 3D zero is an exact focal
cancellation at a positive physical helix height $Z$.  Its chart ordinate is
not $Z$, but
\[
 y=\log Z,
 \qquad
 \texttt{carrierPointAtHeight}(Z)=\sigma_*+i\log Z.
\]
Third, identifying every analytic zero with one of these carrier events is a
separate zero-set coverage statement.

The two proof bodies below are retained as two conditional routes under two
readings: the identity reading and the correlation reading.  The correction is
an exact account of what each displayed premise supplies.  The identity route
assumes analytic-zero identity with a native event; the correlation route
assumes a pointwise chart-zero-to-kernel correspondence.  Neither premise is
now described as a consequence of the internal evidence dossier alone.

This separation is now enforced in the source.  The internal ordinate chart
\texttt{carrierPoint y} uses \texttt{carrierAbscissa}, not a literal half;
the physical-height interface is \texttt{carrierPointAtHeight Z};
\texttt{carrierAbscissa\_eq\_half} derives its normal form from
\texttt{carrierScaleBalanced\_iff}.  The midpoint-free proposition
\texttt{NativeCarrierCoverage}\((\chi)\) says that the zero of ordinate
\(\Im\rho\) produces a completed focal event at physical height
\(Z=e^{\Im\rho}\).  Finally, the compiled decomposition theorem
\[
\texttt{threeDCrossings\_iff\_grh\_and\_nativeCarrierCoverage}
\]
proves the exact content of the older identity premise:
\[
 \texttt{ThreeD\_crossings\_are\_real\_zeros}(\chi)
 \quad\Longleftrightarrow\quad
 \mathrm{GRH}(\chi)\ \wedge\
 \texttt{NativeCarrierCoverage}(\chi).
\]
Accordingly, that older premise is retained as a legacy conditional interface,
not presented as an independently established bridge from analytic zeros to
carrier events.

\medskip
\noindent\textbf{Operator audit.}  At physical height $Z>0$, put
\(\gamma=y=\log Z\).  The local fibre operator used by the elementary 3D
calculation is
\[
 \texttt{vonNeumannOp}(\gamma)=\gamma\,\mathrm{id}_{\mathbb C}.
\]
Its symmetry is a valid theorem, but it is a theorem about a one-dimensional
scalar fibre.  Likewise,
\(\texttt{specBchan}(\gamma,s)=\gamma+i(s-\sigma_*)\) vanishes precisely when
\(s=\sigma_*+i\gamma\).  These facts prove that a point already coupled to
that fibre has the carrier's derived abscissa.  They do not prove that the
nontrivial zeros of \(L(s,\chi)\) exhaust, or are exhausted by, those fibre
events.  This paper therefore calls the scalar operator a \emph{local carrier
coordinate detector}, not a global Hilbert--P\'olya operator.

There is now also one fixed operator on the full geometric 3D carrier,
defined before a character or event predicate is introduced.  For fixed
nondegenerate carrier parameters $(p,r)$ with $p\ne0$, Lean defines
\[
 \mathcal C_{p,r}=\{\,(X,y):X=\gamma_{p,r}(y)\,\},\qquad
 \mathcal H^{(3)}_{p,r}=\mathcal C_{p,r}\to_0\mathbb C,
\]
and
\[
 A^{(3)}_{p,r}=\operatorname{diag}_{s\in\mathcal C_{p,r}}(y_s)
 \quad
 (\texttt{carrierThreeDOperator}).
\]
The state type \texttt{CarrierState3D} stores the spatial point $X$, its
logarithmic ordinate $y$, and only the geometric equation
$X=\texttt{gammaY}(p,r,y)$.  It contains no character, $L$-function, zero
predicate, focal cancellation, or completed-event certificate.  Lean proves
that the third coordinate is $e^y>0$, that \(A^{(3)}_{p,r}\) is symmetric
(\texttt{carrierThreeDOperator\_isSymmetric}), and, for every geometric
carrier state $s$, that \(\delta_s\ne0\) satisfies
\[
 A^{(3)}_{p,r}\delta_s=y_s\delta_s
 \quad
 (\texttt{carrierThreeDOperator\_eigenvector}).
\]
A completed event $e$ is afterward mapped by
\texttt{CompletedThreeDEigenEvent.toCarrierState} to the ambient state with
$y_s=\log Z_e$.  Lean proves that this state's third spatial coordinate is
exactly the original $Z_e$ and that its basis vector is an eigenvector of the
operator already defined on all carrier states
(\texttt{carrierThreeDOperator\_completedEvent\_eigenvector}).  Thus the
event certificate selects a state; it is not used to construct the operator
or its domain.  The older \texttt{completedThreeDOperator} is retained as the
restriction indexed only by completed events for API compatibility.

The Gram pencils also carry an explicit three-coordinate compatibility layer.
For either two-channel Gram matrix $G$, Lean defines
\[
 \texttt{spatialLiftGram}(G)=G\otimes I_3
 \quad\text{on}\quad \mathbb C^2\otimes\mathbb C^3
\]
and proves
\[
 \det(\texttt{spatialLiftGram}(G))=(\det G)^3,
 \qquad
 \det(\texttt{spatialLiftGram}(G))=0
 \Longleftrightarrow \det G=0.
\]
The physical-height definitions \texttt{completedHarmonicGram3DAtHeight} and
\texttt{vonNeumannGram3DAtHeight} instantiate this lift.  The theorem
\[
 \texttt{completedThreeDZeroAtHeight\_twoGram3D\_rankDrop}
\]
sends every completed 3D focal event to both lifted rank drops.  Thus adjoining the three spatial carrier coordinates
cannot create a new vanishing or erase an existing channel cancellation.  This
establishes 3D compatibility of the local pencils.  The ambient carrier
operator above then reads their marked crossings inside a state space that
already existed; the spatial lifts and the event embedding perform distinct
jobs.

\medskip
\noindent\textbf{Zero-set interface.}  The completed native event is
\(\texttt{CompletedThreeDZeroAtHeight}\ \chi\ Z\), definitionally positivity
of the physical height together with exact vanishing of the completed focal
residual.  Lean proves
\[
\begin{split}
 &\texttt{CompletedThreeDZeroAtHeight}\ \chi\ Z\\
 &\qquad\Longleftrightarrow
 L\!\left(\texttt{carrierPointAtHeight}(Z),\chi\right)=0,
 \qquad Z>0,
\end{split}
\]
by \texttt{completedThreeDZeroAtHeight\_iff\_L\_zero}.  Thus a completed 3D
crossing is exactly a genuine zero of the represented $L$-function at its
own readout point.  The separate finite object
\(\texttt{ThreeDZero}\ \chi\ \texttt{bank}\ y\) is the truncated
ordinate-level focal detector and must not be substituted for the completed
event without its limiting bridge.  An arbitrary analytic event is membership
\(\rho\in\texttt{NontrivialZeros}(\chi)\), which includes
\(L(\rho,\chi)=0\) in the critical strip.  The coverage step asks when that
analytic fiber projects to the readout of a physical-height event.

For the principal channel this coverage is now decomposed by a compiled
zero-to-source transfer rather than left as a bare identification.  The
constructor \path{principalZeroAnalyticFiber3D} sends every genuine
\(\rho\in\texttt{ZD.NontrivialZeros}\) to the literal eta-configured
three-dimensional \(\texttt{Vec3}\) fiber at \(\rho\), proves that its
rescaled plane partial sums converge to \(\zeta(\rho)=0\), and supplies the
event-independent ambient helix eigenstate at ordinate \(\Im\rho\) and
physical height \(e^{\Im\rho}\).  This first transfer does not mention a
midpoint.  Before any equality between the analytic parameter and a carrier
readout is requested, the constructor
\path{principalZero_focalCancellation_on_carrier} registers the analytic
closure as a \path{PrincipalZeroNativeCarrierEvent3D}.  The event contains a
typed \path{CarrierState3D} on the Archimedean helix at that same physical
height.  Its theorems \path{.on_archimedeanHelix},
\path{.physicalHeight_eq_focalHeight}, and \path{.ambient_eigenvector} prove
literal helix membership, equality of marked heights, and the ambient 3D
eigenvector equation.  Its \path{.noRadialDrift} theorem follows directly
from helix membership through
\path{carrierPointAtHeight_noRadialDrift}; neither a Euclidean fiber radius
nor a radial-limit premise is used.  The area law is applied only afterwards
by \path{.carrierReadout_scaleBalanced}.  The constructor
\path{principalZeroNative3DSourceTransfer_of_noRadialDrift} then performs
the geometric projection into a conservative source and the complete native
focal/operator certificate from the area-law equation
\(n^{\Re\rho-\sigma_*}=1\).  Its converse is proved from source conservation,
giving the exact theorem
\path{nonempty_principalZeroNative3DSourceTransfer_iff_noRadialDrift}.
Thus the native landing is not encoded by writing
\(\rho=\tfrac12+i z\); the half-unit is the evaluated value of the
area-law-selected \(\sigma_*\) after no drift has been established.

Here ``fiber rides carrier, carrier rides helix'' is a typed composition.
The first map is
\path{PrincipalZeroAnalyticFiber3D.RidesCarrier}, equality of the literal
\(\mathbb N\to\texttt{Vec3}\) analytic fiber with the spectral fiber evaluated
at its carrier-height readout.  The second is the independent geometric field
\path{CarrierState3D.point_eq_gammaY}.  Lean proves that the first map is
equivalent to the parameter-preserving weld by
\path{ridesCarrier_iff_carrierWeld}, and every completed native transfer proves
both maps simultaneously in
\path{fiberRidesCarrier_carrierRidesHelix}.  Merely sharing an ordinate is not
used as a substitute for the fiber attachment.

The exact completed focal event has carrier support by construction and
theorem, not by a later chart assumption.  Namely,
\path{CompletedFocalCancellationRepresents}$(\chi,\rho)$ requires a positive
physical height \(Z\), exact vanishing of \(\texttt{Dcell}\,\chi\,Z\), and the
same-parameter readout \(\texttt{carrierPointAtHeight}\,Z=\rho\).  Lean proves
\path{completedFocalCancellationRepresents_iff}:
\[
 \texttt{CompletedFocalCancellationRepresents}(\chi,\rho)
 \Longleftrightarrow L(\rho,\chi)=0\ \wedge\ \Re\rho=\sigma_*.
\]
Thus \path{no_offCarrier_completedFocalCancellation} rules out an
off-carrier focal cancellation pointwise, before the \(S(t)\) ledger or a
resolvent trace is formed.

The function \(S(t)\) is the global coordinate map for this readout, not a
local numerical coincidence.  In the compiled registration law,
\[
 N_{\pi/3}(e^t)-N_1(e^t)=S(t),
\]
and \texttt{count\_hasJump\_iff\_S\_hasJump} identifies its jumps with the
global analytic event ledger, including multiplicity through
\texttt{residue\_eq\_Smult\_jump}.  Thus \(S(t)\) globally transports the
carrier-scale census to the unit-chart census.  This establishes the global
coordinate bookkeeping used in the correlation route.  Pointwise, the
compiled transfer theorem above identifies native landing with no drift;
globally, \path{allNative3DSourceTransfer_of_selfAdjointXiReceiver} selects that
landing for every nontrivial zero in both half-planes.
Moreover the contour defect cannot hide an off-carrier contribution:
\path{classicalSContour_sub_Smult_pos_of_offCarrier} proves that any enclosed
off-carrier zero makes
\(\texttt{classicalSContour}-\texttt{Smult}>0\), with lower bound its positive
analytic order.  At the operator level,
\path{zeroPoleResolvent_eq_carrierBasisReadout_iff_scaleBalanced} proves
basiswise that the analytic pole receiver equals the independently defined
event-free carrier receiver exactly at area-law balance; inversion is
injective, so neither a trace sum nor a reflected partner can create this
equality.  The corresponding global equivalence is
\path{globalCoordinateIdentification_iff_eachZeroPoleResolvent_eq_carrierBasisReadout}.
The summed bridge is now derived only after this pointwise comparison.  The
definitions \path{xiZeroResolventKernel} and
\path{xiCarrierHeightResolventKernel} retain the zero basis index, analytic
multiplicity, and Hadamard counterterm, while
\path{xiZeroResolventKernel_eq_carrierHeightKernel_iff_scaleBalanced} proves
their equality exactly at area-law balance.  Lean then compiles
\path{globalCoordinateIdentification_iff_basisResolvedCarrierResolventBridge};
after its basiswise equality,
\path{carrierHeightResolventTrace_of_globalCoordinateIdentification} rewrites
the completed xi-channel as the regularized trace of the real carrier-height
receiver.  No equality of summed traces is used to manufacture the weld of an
individual analytic mode.
At the state level,
\path{nonempty_principalZeroNative3DSourceTransfer_iff_fiberRidesCarrier}
proves for each zero that a completed native source exists exactly when the
literal analytic \(\mathbb N\to\texttt{Vec3}\) fiber equals the carrier-readout
fiber at its own physical focal height.  Its quantified form
\path{allNative3DSourceTransfer_iff_allCanonicalZeroFibersRideCarrier} and the
operator form
\path{allNative3DSourceTransfer_iff_zeroLedger_eq_carrierHeight} make the
source/state/operator weld explicit, rather than treating the common
ordinate as sufficient.
The operator propagation is then literal on the analytic-zero basis:
\path{xiZeroLedgerOperator_eq_carrierHeight_of_allNative3DSourceTransfer}
uses each completed transfer to prove that its zero-ledger coefficient equals
its real carrier-height coefficient.  Consequently
\path{globalCoordinateIdentification_of_allNative3DSourceTransfer} returns
\(\texttt{classicalSContour}=\texttt{Smult}\) at every good height, without a
summed-trace cancellation step.  The receiver capstone
\path{globalCoordinateIdentification_and_zeroLedger_eq_carrierHeight_of_selfAdjointXiReceiver}
returns the global coordinate identity and the operator equality together.
Under the stated (Z--3D) reading that a native 3D zero is primary and its 1D
analytic zero is its shadow, this assembly is now exposed by one theorem:
\path{globalCoordinateIdentification_basisResolvedBridge_and_carrierTrace_of_allNative3DSourceTransfer}.
Its sole conditional input is one completed native source transfer for every
analytic shadow.  It returns simultaneously the global coordinate identity,
the equality of the two resolvent kernels on every zero basis mode, and the
regularized carrier-height trace.  Thus the shadow conditional is used once
at the source/state weld; no equation \(\Re\rho=1/2\) is supplied as an
additional premise to the capstone.

\medskip
\noindent\textbf{Formal dependency table.}
\begin{center}
\small
\renewcommand{\arraystretch}{1.15}
\begin{tabular}{p{0.31\linewidth}|p{0.22\linewidth}|p{0.35\linewidth}}
Statement & Proven here & Logical role\\
\hline
\(\texttt{carrierAbscissa}=\tfrac12\) & yes, from the area law & fixes the
native carrier image\\
\shortstack[l]{\texttt{PrincipalZero}\\
\texttt{AnalyticFiber3D}\((\rho)\)} & yes, for every genuine
zeta zero & literal 3D analytic state, zero readout, physical height and
ambient eigenvector\\
\shortstack[l]{\texttt{PrincipalZero}\\
\texttt{NativeCarrierEvent3D}\((\rho)\)} & yes, for every genuine
zeta zero & registers focal closure on a typed helix state at the same height;
no drift follows from carrier membership, with no radius comparison\\
\shortstack[l]{\texttt{PrincipalZero}\\
\texttt{Native3DSourceTransfer}\((\rho)\)} & yes, from no drift; converse proved &
connects the analytic fiber to its conservative source and complete focal
certificate without a midpoint input\\
\shortstack[l]{\texttt{ThreeD\_crossings\_are}\\
\texttt{real\_zeros}\((\chi)\)} & legacy premise &
equals \(\mathrm{GRH}(\chi)\wedge\texttt{NativeCarrierCoverage}(\chi)\)\\
\shortstack[l]{\texttt{OneD\_zeta\_zero}\\
\texttt{correlated}\((\chi)\)} & yes, by its right branch &
routing only; it supplies no kernel identification\\
\shortstack[l]{\texttt{OneDCorrelation}\\
\texttt{Evidence}\((\chi)\)} & a separate premise & places a
scalar-fibre kernel at every analytic chart zero and by itself implies the
critical-line conclusion\\
\end{tabular}
\end{center}

\medskip
\noindent\textbf{What the supporting facts establish.}  The finite phasor,
Gram, projection, counting, and convergence theorems certify the internal
carrier mechanics and several chart-side audits.  A diagonal operator indexed
by already located analytic zeros checks counts and multiplicities; it does
not locate those zeros.  Similarly, local identities between a focal residual
and an \(L\)-value must be read with their stated domains and limiting
hypotheses.  None of these facts is used below as a substitute for the
analytic-zero-to-native-event bridge.

\medskip
\noindent\textbf{Scope.}  The midpoint theorem is independent of
\(L(s,\chi)\): it is a geometric statement about which fibre scaling remains
finite and drift-free on the Archimedean helix.  The GRH conditionals are
different statements about whether analytic zeros are represented on that
geometric image.  The revised Lean declarations and theorems expose this
boundary rather than asking the reader to infer it from prose.

\section{The carrier construction and its interfaces}\label{sec:grh-construction}

Everything in this section is Part~I machinery specialized to the basic
configuration; every object carries a pointer to its general construction.
The objects come first --- the state space, the carrier that lives in it, the
fiber that rides the carrier --- then the maps between scales, then the
events, then the operators that read them.

\subsection{The 3D state space}\label{sec:grh-statespace}

The ambient space is three-dimensional, with named coordinates: \emph{radius}
\(r\), \emph{phase} \(\theta\), and physical \emph{height} \(Z\). A complete carrier
coordinate state is a point
\[
\mathsf X=(r,\theta,Z)\in\mathbb R_{\ge0}\times\mathbb R\times\mathbb R_{\ge0},
\]
and this is the space in which everything below lives: the carrier is a curve
in it, a phasor is a vector in the transverse \((r,\theta)\)-plane at its
height, and the \emph{bank state} at height \(Z\) is the accumulated sum of
every phasor that has entered below \(Z\) --- one such state per lane. The 3D
representation of \(L(s,\chi)\) is exactly this pair of lane states, evolving
continuously as \(Z\) grows; the operators of \S\ref{sec:grh-operators} read
it, and the events of \S\ref{sec:grh-focal} are its alignments.

Height is the distinguished coordinate. It is the one coordinate the 1D chart
retains (\S\ref{sec:grh-projection}), every event in this part is indexed by
the height at which it occurs, and there is no convergence gate anywhere in
the space: states exist and evolve at every height, and the half-plane
\(\Re s>1\) of the classical series is a property of the compressed chart,
not of the state space.

\subsection{The carrier: a double-ended helix on the rescaled number
line}\label{sec:grh-helix}

The carrier is the source-independent geometric object in the state space ---
fixed, by Theorem~\ref{thm:g1-double-carrier}, before any character is
attached. It is double-ended: two lanes \(C=C_+\sqcup C_-\) over a common raw
parameter, realized as
\[
\Gamma_{\pm}(t)=\bigl(R(t)\cos(\pm\Theta(t)),\,R(t)\sin(\pm\Theta(t)),\,Z(t)\bigr),
\]
with the global involution \(J:C_+\leftrightarrow C_-\), \(J^2=\mathrm{id}\),
exchanging the lanes while preserving radius and height and reversing the
transverse phase. The plus lane is the helix; the minus lane is its conjugate
anti-helix at the same height. In the basic configuration the helix is
Archimedean with constant pitch of one unit --- one unit of height per turn,
the transverse shadow an Archimedean spiral --- so the winding, radius, and
height laws grow together with no free parameter.

The carrier admits an elementary description: it is the positive number line,
rescaled. The height coordinate \emph{is} the number line --- the integer
\(n\) sits at height \(Z=n\) --- and the rescaling is harmonic: a unit step
of the integers advances the winding by \(\pi/3\). Six steps close one full
turn, so the finite harmonic cells are the \(\mu_6\) buckets, and the natural
coefficient ring of the closure arithmetic is the Eisenstein ring
\(\mathbb Z[\zeta_6]\), in which the cancellation at a closed cell is exact
--- integer arithmetic, no limit taken (\S\ref{sec:construct}). Nothing about
any source is encoded in this choice; it is the harmonic scale at which the
carrier's cells close exactly, and every integer occupies the same carrier
whether or not any character ever reads it.

The two lanes are not decorative. The functional equation of the readout is
the involution \(J\) read phasor by phasor: on the weld axis, height inversion
coincides with conjugation, and \(J\) supplies exactly that exchange
(Theorem~\ref{thm:carrierreflection}, Lemma~\ref{lem:bankreflection}). What
the classical theory obtains from Poisson summation, the carrier carries as a
geometric symmetry of the state space, fixed before any source is attached.
The weld axis --- the fixed locus of \(J\) --- is the carrier's copy of
Riemann's real axis: the midline at \(\textsc{unit}/2\), derived from the
arclength area law and no-drift conservation in
\S\ref{sec:grh-midpoint-provenance}, then expressed in the unit-width chart
of \S\ref{sec:grh-overview}. Figure~\ref{fig:grh-weld} shows where the origin
and the functional equation enter: the double-ended object welded at
\(z=1\), the height inversion pivoting there, and \(J\) supplying it as
conjugation on the axis.

\begin{figure}[t]
\centering
\resizebox{0.82\textwidth}{!}{%
\begin{tikzpicture}[>={Stealth[length=2.4mm]}, line cap=round]
\colorlet{carrierY}{yellow!85!orange}
\colorlet{fiberB}{blue!72!black}
\colorlet{midV}{violet!85!blue}
\colorlet{vanishR}{red!82!black}
% ---- background ----
\fill[blue!4, rounded corners=8pt] (-6.6,-5.3) rectangle (6.6,5.6);
\node[anchor=north west, font=\bfseries\small, text=blue!40!black]
  at (-6.45,5.50) {the double-ended object: origin and functional equation};
% ---- weld axis (midline) ----
\draw[midV, line width=1.4pt, densely dash dot] (0,-4.35) -- (0,4.35);
% ---- weld plane at z=1 ----
\fill[gray!16, opacity=0.75] (0,0) ellipse [x radius=2.35, y radius=0.42];
\draw[gray!60, densely dashed] (0,0) ellipse [x radius=2.35, y radius=0.42];
\node[anchor=west, font=\scriptsize, align=left] at (2.55,0.02)
  {weld plane: $z=1$ ($y=0$;\\[-2pt]$x=1$ in the theta variable)};
% ---- upper helix: readout leg ----
\draw[carrierY, line width=4.2pt, opacity=0.95]
  plot[domain=0:1800, samples=320, smooth]
  ({(0.55+0.0003*\x)*cos(\x)}, {\x/450+0.2*(0.55+0.0003*\x)*sin(\x)});
\draw[fiberB, line width=1.25pt]
  plot[domain=0:1800, samples=320, smooth]
  ({(0.55+0.0003*\x)*cos(\x)}, {\x/450+0.2*(0.55+0.0003*\x)*sin(\x)});
% ---- lower anti-helix: dual leg (exact mirror through the weld) ----
\draw[carrierY, line width=4.2pt, opacity=0.95]
  plot[domain=0:1800, samples=320, smooth]
  ({(0.55+0.0003*\x)*cos(\x)}, {-\x/450-0.2*(0.55+0.0003*\x)*sin(\x)});
\draw[fiberB, line width=1.25pt]
  plot[domain=0:1800, samples=320, smooth]
  ({(0.55+0.0003*\x)*cos(\x)}, {-\x/450-0.2*(0.55+0.0003*\x)*sin(\x)});
% ---- origin / weld point ----
\fill[midV] (0,0) circle (2.8pt);
\node[midV, anchor=north west, font=\scriptsize\bfseries] at (0.14,-0.30)
  {the origin/weld --- central point $s=\tfrac12$};
% ---- continuations to infinity ----
\draw[fiberB, densely dashed, ->] (1.09,4.02) -- (1.32,4.62)
  node[anchor=west, font=\scriptsize] {to $\infty$};
\draw[fiberB, densely dashed, ->] (1.09,-4.02) -- (1.32,-4.62)
  node[anchor=west, font=\scriptsize] {the $z<1$ end, folded};
% ---- functional-equation reflection P <-> P^vee ----
\fill[fiberB, opacity=0.18] (0.874,2.4) circle (0.22);
\fill[fiberB, opacity=0.18] (0.874,-2.4) circle (0.22);
\fill[fiberB] (0.874,2.4) circle (2.4pt);
\fill[fiberB] (0.874,-2.4) circle (2.4pt);
\node[anchor=west, font=\scriptsize] at (1.10,2.58) {$s=\tfrac12+iy$};
\node[anchor=west, font=\scriptsize] at (1.10,-2.58) {$1-s=\tfrac12-iy$};
\draw[<->, vanishR, thick, densely dashed]
  (1.05,2.32) .. controls (3.9,1.3) and (3.9,-1.3) .. (1.05,-2.32);
\node[vanishR, align=center, font=\scriptsize, anchor=west] at (3.45,0.95)
  {height inversion\\[-2pt]$z\mapsto 1/z$ \ ($s\mapsto 1-s$)};
% ---- J and epsilon annotations (left) ----
\node[align=right, anchor=east, font=\scriptsize] at (-2.55,1.05)
  {on the weld axis, inversion $=$ conjugation:\\[-2pt]
   $J$, phasor by phasor};
\draw[densely dotted] (-2.45,1.05) -- (-0.12,0.12);
\node[align=right, anchor=east, font=\scriptsize] at (-2.55,-1.30)
  {$\varepsilon=\prod_v(-\bar\alpha_v)$, $|\varepsilon|=1$:\\[-2pt]
   the offset at the weld};
\draw[densely dotted] (-2.45,-1.30) -- (-0.12,-0.14);
% ---- leg labels ----
\node[align=right, anchor=east, font=\scriptsize, fiberB] at (-2.55,3.10)
  {readout leg $\Lambda(s;W)$:\\[-2pt]heights $z>1$, fiber $W$};
\node[align=right, anchor=east, font=\scriptsize, fiberB] at (-2.55,-3.10)
  {dual leg $\varepsilon\,\Lambda(1-s;W^{\vee})$:\\[-2pt]
   conjugate phases, anti-helix};
% ---- FE banner ----
\node[font=\scriptsize\bfseries] at (0,-4.95)
  {$\Lambda(s;W)=\varepsilon\,\Lambda(1-s;W^{\vee})$ --- the strand
   exchange, read at the weld};
\end{tikzpicture}}
\caption{Where the origin and the functional equation enter. The carrier is
double-ended: the readout leg (heights \(z>1\), fiber \(W\)) and the dual
leg (the \(z<1\) end, conjugate phases on the anti-helix, fiber
\(W^{\vee}\)) are welded at the origin \(z=1\) --- the central point
\(s=\tfrac12\), where the weld plane meets the weld axis. The functional
equation is the fold through that weld: height inversion \(z\mapsto1/z\)
(\(s\mapsto1-s\)) exchanges the legs, and on the weld axis the inversion
coincides with conjugation, which \(J\) supplies phasor by phasor
(Theorem~\ref{thm:carrierreflection}, Lemma~\ref{lem:bankreflection},
step~1) --- no Poisson summation is consumed. The only mismatch the exchange
permits is the unimodular offset at the weld,
\(\varepsilon=\prod_v(-\bar\alpha_v)\), assembled per place in Lean
(\texttt{StrandExchange.bankProduct\_exchange} from
\texttt{ChiralityHB.symClock\_star}); \(\varepsilon\) is the root number,
the offset that places or doubles the central cancellation.}
\label{fig:grh-weld}
\end{figure}

\subsection{The fiber: the character's occupation of the
carrier}\label{sec:grh-buckets}

The carrier is common; the \emph{fiber} is what a particular source puts on
it. A Dirichlet character \(\chi\bmod q\) attaches its fiber \(F_\chi\): one
phasor per integer per lane, with magnitude and phase synthesized from
\(\chi\)'s finite data alone
(Theorem~\ref{thm:g3-g4-identity-synthesis}) --- the phasor of \(n\) carries
the character phase \(\chi(n)\) on the carrier's winding at \(n\). The
carrier \emph{hosts} the fiber; it does not become it: the same carrier point
under two characters carries two different fiber states, distinguished by the
retained identity, and neither a shared carrier point nor a shared scalar
readout identifies two sources
(Corollary~\ref{cor:g3-complete-state-identity}). The 3D representation of
\(L(s,\chi)\) referred to throughout this part is the \emph{realized} fiber:
the complete state of Part~I --- carrier, fiber, native clock, native cells
--- specialized to the basic configuration.

Each phasor enters at its own height and at zero magnitude: the phasor of
\(n\) is born at height \(Z=n\) and grows continuously with height ---
there is no convergence gate on the carrier
(\S\ref{sec:grh-statespace}). The dual copy rides the anti-helix with the
same height law and conjugate transverse phase, so the two lanes run in
parallel, heights matching, geometries conjugate
(Theorem~\ref{thm:g1-double-carrier}).

The conductor decides the buckets. For \(n\) sharing a factor with \(q\) we
have \(\chi(n)=0\): the phasor is \emph{neutral} (U) --- present on the
carrier, silent in this character's readout. For \(n\) coprime to \(q\) the
character phase routes the phasor: a real character splits the units into a
\emph{positive} lane P (\(\chi(n)=+1\)) and a \emph{negative} lane M
(\(\chi(n)=-1\)); a complex character distributes its phases over the
\(\mu_6\)-cell directions, and the P/M split is realized as the signed channel
pair \((A,B)\) --- the \(\operatorname{Re}\)- and \(\operatorname{Im}\)-lanes
of the reflected readout (Theorem~\ref{thm:g6-mode-bridge}). The same integers
occupy the same carrier for every character; only the bucket assignment
changes. 

\subsection{The bidirectional transport: down to the ordinate, back to the
state}\label{sec:grh-projection}

The transport between the 3D state and the 1D chart runs in both directions,
each direction a named map, and both composites are the identity
(Theorem~\ref{thm:g2-ledgered-projection}). Neither direction is a metaphor:
down is how the classical function is read off the state, and up is how a
chart event is located back on the carrier.

\emph{Downward.} The descent is two coordinate suppressions and a logarithm.
Write a complete carrier coordinate state as \(\mathsf X=(r,\theta,z)\)
(\S\ref{sec:grh-statespace}). The first projection forgets the radius,
\[
\Pi_{32}(r,\theta,z)=(\theta,z),
\]
recording the two-dimensional carrier image --- in the reflected chart, the
Cayley/M\"obius image on the unit-circle chart
(Theorem~\ref{thm:carrierreflection}, item (P)). The second forgets the
phase,
\[
\Pi_{21}(\theta,z)=z,
\]
leaving the height alone. The analytic chart then reads the retained height
logarithmically,
\[
\Pi^{\log}_{21}(\theta,z)=\log z=y,
\]
and places it at the represented point \(\tfrac12+iy\) of the strip: the
abscissa supplied by the derived weld axis
(\S\S\ref{sec:grh-helix} and~\ref{sec:grh-midpoint-provenance}), the
ordinate by the logarithm of the height. That is
the whole downward mechanism. A three-dimensional value becomes a
one-dimensional ordinate by having its radius suppressed, its phase
suppressed, and its height passed through \(\log\). What is suppressed is
booked, not destroyed: each projection is paired with the loss ledger
\(\mathcal D(\mathsf X)=(r,\theta)\) --- in general: radius, angle, clock
offset, adapter state, every coordinate the chosen chart makes invisible.

\emph{Upward.} The ascent is reconstruction from the readout and the ledger,
and it is exact:
\[
\operatorname{Rec}\bigl(\Pi(\mathsf X),\mathcal D(\mathsf X)\bigr)=\mathsf X,
\qquad
\operatorname{Rec}^{\log}\bigl(y,(r,\theta)\bigr)=(r,\theta,e^{y}).
\]
A chart ordinate \(y\) lifts to the carrier height \(Z=e^{y}\); the ledger
restores the radius and phase; and the lift lands on the exact state the
projection left. This is a bijection with an explicit two-sided inverse ---
\(\texttt{record}\,f=(\text{height},(\text{radial},\text{phase}))\) inverts
by \texttt{reconstruct} and \texttt{reconstruct} inverts by \texttt{record}
(Lean \texttt{ConeProjection.record\_bijective}, \texttt{reconstruct\_record});
distinct helix events write distinct ledger entries
(\texttt{SourceHolonomy.projection\_bijective\_loss\_ledger}). Transport with
the ledger is lossless in both directions; transport without the ledger is a
compression, and provably so --- dropping the radial channel alone already
makes the projection non-injective (\texttt{ConeProjection.radial\_lost},
\S\ref{sec:grh-E1}).

The two directions do different work in this part. Downward is
\emph{reading}: the value of \(L(s,\chi)\) at the represented point
\(\tfrac12+iy\) is the final verification of a marked event --- the analytic
function is consulted after the event, never used to build it
(non-circularity, \S\ref{sec:grh-overview}). Upward is \emph{locating}: a
chart ordinate \(y\) names the carrier height \(e^{y}\) at which the operator
reads the bank, and the two reading-hypotheses of this part are statements
about precisely this lift --- the 1D-primary branch of the weak hypothesis is
the upward lift \(\gamma=\Im\rho\) landing on an eigenstate. The classical
analytic function is the downward readout \emph{without} the ledger --- the
ledger stays on the carrier. This is why equality of readouts does not imply
equality of states: two states share an ordinate whenever their heights
agree, whatever their radii and phases are doing; the exact criterion is that
readout \emph{plus ledger} determines the state
(Corollary~\ref{cor:g2-projected-equality}).

Both decisions of \S\ref{sec:grh-overview} take their teeth from this
mechanism. The 1D ordinate is the logged trace of the 3D event; the
eigenstate carries the coordinates the ordinate dropped. Whether an
eigenvalue must \emph{be} that trace or \emph{coincide} with it (decision
one), and whether the event or its trace is ``the zero'' (decision two), are
questions about the two ends of the maps displayed above.
Figure~\ref{fig:grh-pipeline} draws the entire pipeline for \(\chi_3\): the
conjugate pair up to its first crossing, and the ledgered descent to the
chart --- with the midline surviving every stage.

\begin{figure}[p]
\centering
\resizebox{\textwidth}{!}{%
\begin{tikzpicture}[>={Stealth[length=2.4mm]}, line cap=round]
\colorlet{carrierY}{yellow!85!orange}
\colorlet{fiberB}{blue!72!black}
\colorlet{midV}{violet!85!blue}
\colorlet{vanishR}{red!82!black}
% ================= PANEL BACKGROUNDS =================
\fill[blue!4, rounded corners=8pt] (-8.0,-1.75) rectangle (4.35,8.85);
\fill[gray!7, rounded corners=8pt] (4.65,-1.75) rectangle (15.4,8.85);
\node[anchor=north west, font=\bfseries\small, text=blue!40!black]
  at (-7.85,8.75) {3D state space};
\node[anchor=north east, font=\bfseries\small, text=black!60]
  at (15.25,8.75) {the ledgered descent};
% ================= PANEL A: the 3D object =================
% weld axis: the midline, unit/2 (midpoint no. 1)
\draw[midV, line width=1.4pt, densely dash dot] (0,-0.15) -- (0,7.55);
\node[midV, font=\scriptsize\bfseries] at (0,-1.30)
  {midpoint: weld axis $=\textsc{unit}/2$};
% height axis
\draw[->, thick] (-5.15,0) -- (-5.15,7.9) node[above] {\small height $Z$};
\foreach \n in {1,2,3}{
  \draw (-5.23,\n) -- (-5.07,\n) node[left=3pt] {\scriptsize $n=\n$};}
\draw[thick] (-5.29,5.25) -- (-5.01,5.45);
\draw[thick] (-5.29,5.10) -- (-5.01,5.30);
\draw (-5.23,7) -- (-5.07,7);
\node[anchor=east, font=\scriptsize] at (-5.30,7)
  {$Z=e^{y}\approx 3.10\times10^{3}$};
% ---- helix C+ (left): yellow carrier band, blue fiber riding it ----
\draw[carrierY, line width=4.2pt, opacity=0.95]
  plot[domain=0:3420, samples=480, smooth]
  ({-1.7+(0.55+0.00017*\x)*cos(\x)}, {\x/488.57+0.22*(0.55+0.00017*\x)*sin(\x)});
\draw[fiberB, line width=1.25pt]
  plot[domain=0:3420, samples=480, smooth]
  ({-1.7+(0.55+0.00017*\x)*cos(\x)}, {\x/488.57+0.22*(0.55+0.00017*\x)*sin(\x)});
% ---- anti-helix C- (right, mirrored winding, conjugate tilt) ----
\draw[carrierY, line width=4.2pt, opacity=0.95]
  plot[domain=0:3420, samples=480, smooth]
  ({1.7-(0.55+0.00017*\x)*cos(\x)}, {\x/488.57-0.22*(0.55+0.00017*\x)*sin(\x)});
\draw[fiberB, line width=1.25pt]
  plot[domain=0:3420, samples=480, smooth]
  ({1.7-(0.55+0.00017*\x)*cos(\x)}, {\x/488.57-0.22*(0.55+0.00017*\x)*sin(\x)});
% ---- phasor entries on C+ ----
\draw[->, fiberB, thick] (-2.34,1.10) -- ++(-0.26,0);
\draw[->, fiberB, thick] (-0.97,2.21) -- ++(0.30,0);
\draw[->, fiberB, thick] (-2.53,3.32) -- ++(-0.36,0);
% ---- left label stack ----
\node[align=left, anchor=west, font=\scriptsize] at (-7.85,4.90)
  {carrier: $\pi/3$-scaled\\[-2pt](yellow layer)};
\draw[densely dotted] (-6.05,4.90) -- (-2.72,4.62);
\node[align=left, anchor=west, font=\scriptsize, fiberB] at (-7.85,3.70)
  {fiber $F_{\chi_3}$ (blue)\\[-2pt]riding the carrier};
\draw[fiberB, densely dotted] (-6.00,3.70) -- (-2.60,3.45);
\node[fiberB, align=left, anchor=west, font=\scriptsize] at (-7.85,2.20)
  {phasor of $n$ enters\\[-2pt]at height $n$, grows};
\draw[fiberB, densely dotted] (-6.00,2.20) -- (-2.66,1.12);
% ---- crossing / focal cancellation (glow) ----
\fill[vanishR, opacity=0.12] (0,7) circle (0.52);
\fill[vanishR, opacity=0.22] (0,7) circle (0.30);
\draw[vanishR, densely dashed, line width=0.9pt] (-3.35,7) -- (3.35,7);
\fill[fiberB, opacity=0.18] (-2.83,7) circle (0.26);
\fill[fiberB, opacity=0.18] (2.83,7) circle (0.26);
\fill[fiberB] (-2.83,7) circle (2.6pt);
\fill[fiberB] (2.83,7) circle (2.6pt);
\node[vanishR] at (0,7) {$\otimes$};
\node[vanishR, align=center, font=\scriptsize, fill=blue!4, inner sep=1.5pt]
  at (0,6.22) {focal cancellation\\[-2pt](vanishing)};
% ---- Frobenius bracket ----
\draw[decorate, decoration={brace, amplitude=7pt}, thick]
  (-2.83,7.42) -- (2.83,7.42);
\node[font=\scriptsize] at (0,8.02)
  {conjugate eigenstates --- Frobenius $\det=1$};
% ---- lane labels ----
\node[font=\scriptsize] at (-2.4,-0.62) {helix $C_{+}$ (phase $+\theta$)};
\node[font=\scriptsize] at (2.4,-0.62) {anti-helix $C_{-}$ (phase $-\theta$)};
% ================= PANEL B: the ledgered descent =================
% ledger box
\node[draw, thick, rounded corners, align=center, inner sep=5pt,
      fill=white, font=\scriptsize] (ledger) at (7.05,6.85)
  {loss ledger $\mathcal D=(r,\theta)$\\[-1pt]radius $r$ booked};
\draw[->, very thick, black!70] (2.95,7.05) .. controls (4.55,7.10) .. (ledger.west);
% ---- 2D stage: unit circle (midpoint no. 2) ----
\draw[midV, line width=1.3pt] (11.6,6.15) circle (0.95);
\fill (11.6,6.15) circle (1.1pt);
\draw[gray, densely dashed] (11.6,6.15) -- ++(128:0.95);
\fill[fiberB, opacity=0.18] ($(11.6,6.15)+(128:0.95)$) circle (0.2);
\fill[fiberB] ($(11.6,6.15)+(128:0.95)$) circle (2.2pt);
\draw (11.95,6.15) arc[start angle=0, end angle=128, radius=0.35];
\node[font=\scriptsize] at (11.85,6.60) {$\theta$};
\node[align=center, font=\scriptsize] at (11.6,4.82)
  {Cayley/M\"obius $\to$ 2D unit circle\\[-2pt](radius held in the ledger)};
\node[midV, font=\scriptsize\bfseries, align=center] at (11.6,7.55)
  {midpoint: lands \emph{on} the circle};
\draw[->, very thick, black!70] (ledger.east) -- (10.55,6.45);
% ---- 1D stage ----
\draw[->, very thick, black!70] (11.6,4.45) -- (11.6,3.85)
  node[midway, right=3pt, font=\scriptsize] {project: angle $\theta$ booked};
\draw[thick] (10.2,3.55) -- (13.0,3.55);
\fill[fiberB] (11.6,3.55) circle (2.2pt);
\node[anchor=west, font=\scriptsize] at (13.1,3.62)
  {1D: height retained};
\node[anchor=west, font=\scriptsize, text=black!60] at (13.1,3.22)
  {abscissa pinned at $\textsc{unit}/2$};
% ---- chart stage (midpoint no. 3) ----
\draw[->, very thick, black!70] (11.6,3.25) -- (11.6,2.62)
  node[midway, right=3pt, font=\scriptsize] {$\log$};
\fill[white] (10.6,0.42) rectangle (12.6,2.32);
\draw[thick] (10.6,0.42) rectangle (12.6,2.32);
\draw[midV, line width=1.3pt, densely dash dot] (11.6,0.42) -- (11.6,2.32);
\fill[fiberB, opacity=0.18] (11.6,1.75) circle (0.2);
\fill[fiberB] (11.6,1.75) circle (2.2pt);
\node[anchor=west, font=\scriptsize] at (12.7,1.75)
  {$y=\log Z\approx 8.04$};
\node[midV, font=\scriptsize\bfseries] at (11.6,0.14)
  {midpoint: $\Re s=\tfrac12$};
\node[font=\scriptsize] at (11.6,-0.42)
  {readout: first zero of $L(s,\chi_3)$};
\node[font=\scriptsize] at (11.6,-0.86)
  {chart registration: $N_{\pi/3}(e^{t})-N_{1}(e^{t})=S(t)$};
% ---- return arrow (bidirectional transport) ----
\draw[->, densely dashed, thick, gray!65!black]
  (10.5,1.35) .. controls (8.0,1.7) and (6.35,3.6) .. (6.70,6.35);
\node[gray!65!black, align=center, rotate=68, font=\scriptsize] at (6.82,3.65)
  {$\operatorname{Rec}$: $Z=e^{y}$; ledger restores $(r,\theta)$};
% ---- midline thread annotation ----
\node[midV, font=\scriptsize\bfseries] at (10.0,-1.38)
  {midpoint $\to$ midpoint $\to$ midpoint: the midline survives every stage};
\end{tikzpicture}}
\caption{The pipeline for \(\chi_3\), end to end (heights schematic, not to
scale). \emph{Left panel} --- the 3D object: helix \(C_{+}\) and conjugate
anti-helix \(C_{-}\), the \(\pi/3\)-scaled carrier as the yellow layer, the
fiber \(F_{\chi_3}\) in blue growing from the origin with phasors entering
at integer heights; the violet dash-dot line is the weld axis --- the
midline at \(\textsc{unit}/2\), the fixed locus of \(J\). At the first
crossing the lanes align in exact focal cancellation, and the conjugate
eigenstate pair (Frobenius \(\det=1\)) marks the event at height
\(Z=e^{y}\approx3.10\times10^{3}\). \emph{Right panel} --- the ledgered
descent: radius booked in the loss ledger, Cayley/M\"obius to the 2D unit
circle, angle booked at the projection to 1D, logarithm to the chart,
landing the readout at \(y=\log Z\approx 8.0397\), the first zero of
\(L(s,\chi_3)\), with \(S(t)\) the registration bookkeeping between the two
scales (\S\ref{sec:carrier-S}). The violet thread is the midline surviving
every stage --- weld axis in 3D, on-circle landing in 2D
(\texttt{midpoint\_entry\_on\_circle}), \(\Re s=\tfrac12\) in the chart
(\texttt{pipeline\_midpoint\_iff}): the ordinate compresses by a logarithm
while the abscissa is first derived by the area law and then held by
no-drift symmetry. The dashed return arrow is the
upward transport: readout plus ledger reconstructs the state exactly
(\S\ref{sec:grh-projection}).}
\label{fig:grh-pipeline}
\end{figure}

% Midpoint audit searches run 2026-07-14:
% rg -n -i "carrierPoint|carrier_point|ThreeD_crossings_are_real_zeros|three_d_crossings_are_real_zeros|crossings.*real.*zeros|ThreeDZero|three_d_zero|no.?drift|drift.?zero|offline|off.?line|off.?carrier|midpoint|half.?unit|area.?law|archimedean.*helix|helix.*archimedean" . --glob '*.lean' --glob '*.tex' --glob '!staging_rp2/**'
% rg -n -i "(no.*drift|drift.*zero|drift.*iff|scaleBalanced_iff|sigma_half_is_scale_critical|carrier.*offline|offline.*carrier|cancellation.*offline|offline.*cancellation|admissible.*offline|offline.*admissible|reprPoint|spectralCoord|carrierPoint|pipeline_midpoint_iff)" RequestProject --glob '*.lean'
% rg -n -i "carrierPoint|carrier_point|ThreeD_crossings_are_real_zeros|three_d_crossings_are_real_zeros|ThreeDZero|three_d_zero|no_radial_drift_on_helix|helix_no_offline_cancellation|sigma_half_is_scale_critical|scaleBalanced_iff" .lake/packages/mathlib/Mathlib --glob '*.lean'

\subsection{Midpoint provenance: geometry first, chart second}
\label{sec:grh-midpoint-provenance}

The order of construction matters here.  The carrier is not introduced as
the vertical line \(\Re s=\tfrac12\).  Before there is an \(s\)-plane, a
zero \(\rho\), or a definition of \texttt{carrierPoint}, the wound integer
line has a cylindrical radius \(r_n\).  The arclength calculation gives
\[
  \frac{r_n^2}{n}\longrightarrow \frac{r\Delta}{\pi}>0,
  \qquad
  \frac{r_n}{\sqrt n}\longrightarrow
     \sqrt{\frac{r\Delta}{\pi}}>0
\]
(\texttt{windIntegerSite\_radius\_sq\_tendsto} and
\texttt{carrierRadius\_div\_sqrt\_tendsto}).  Thus the exponent
\(\sigma\) of a fiber amplitude \(n^{-\sigma}\) is initially free.  Define
geometric admissibility by the existence of a finite positive scale limit,
\[
 \operatorname{ScaleBalanced}(\sigma)
 \quad:\Longleftrightarrow\quad
 \exists L>0,
 \quad n^{-\sigma}r_n\longrightarrow L.
\]
The area law then proves
\[
 \operatorname{ScaleBalanced}(\sigma)
 \quad\Longleftrightarrow\quad
 \sigma=\frac12.
 \tag{MP1}\label{eq:grh-midpoint-area-law}
\]
This is \texttt{sigma\_half\_is\_scale\_critical}, re-exported as
\texttt{scaleBalanced\_iff}.  The gauge parameters \(r\) and \(\Delta\)
change the positive limiting constant, but not the exponent.  Consequently
the midpoint is an output of the three-dimensional arclength geometry: it
is the exponent of the square-root radius, not a coordinate selected from
the analytic zero set.

There is an independent no-drift form of the same conclusion.  A raw
source mode has a complex transport rate \(\lambda\) and nonzero amplitude
\(a\), with conservation law
\[
  e^{(\Re\lambda)\tau}a=a\qquad(\tau\in\mathbb R).
\]
Taking \(\tau=1\) forces \(\Re\lambda=0\)
(\texttt{source\_noDrift} and \texttt{SourceMode.noDrift}).  This statement
contains no \(\rho\), no \(s\), and no numerical midpoint.  In the radial
normalization supplied by (MP1), a putative abscissa \(\sigma\) has residual
drift \(n^{\sigma-\sigma_*}\), where
\(\sigma_*\) is the unique scale-balanced exponent.  Hence
\[
  n^{\sigma-\sigma_*}=1\ (n>1)
  \quad\Longleftrightarrow\quad \sigma=\sigma_*
  \quad\Longleftrightarrow\quad \sigma=\frac12,
 \tag{MP2}\label{eq:grh-midpoint-drift}
\]
the compiled specialization being
\texttt{Helix.no\_radial\_drift\_iff\_half}.  Thus the two inputs agree:
the area law determines the centre \(\sigma_*\), and conservation prohibits
motion away from it.

Only after (MP1)--(MP2) is the spectral chart attached.  The physical height
and its internal logarithmic ordinate are
\[
 Z>0,\qquad y=\log Z,\qquad
 \operatorname{carrierPointAtHeight}(Z)=\sigma_*+i\log Z.
\]
Lean stores the ordinate chart and physical-height readout separately:
\[
\begin{aligned}
 \texttt{carrierPoint y}&:=\texttt{carrierAbscissa}+iy,\\
 \texttt{carrierPointAtHeight Z}&:=\texttt{carrierPoint}(\log Z).
\end{aligned}
\]
Here \texttt{carrierAbscissa} is selected from the unique witness of
\texttt{CarrierScaleBalanced}; only the theorem
\texttt{carrierAbscissa\_eq\_half} rewrites it to \(\tfrac12\).  The theorem
\texttt{carrierPoint\_eq\_half\_add\_I} is therefore a downstream normal-form
lemma, not the definition of the carrier coordinate.  Theorems
\texttt{carrierPointAtHeight\_im} and
\texttt{carrierPointAtHeight\_exp} certify $y=\log Z$ and
$Z=e^y$ in both directions.

The completed native crossing is
\[
\texttt{CompletedThreeDZeroAtHeight}\ \chi\ Z
 \quad:\Longleftrightarrow\quad
 Z>0\ \wedge\ \texttt{Dcell}\ \chi\ Z=0.
\]
The exact representation theorem gives
\[
 \texttt{CompletedThreeDZeroAtHeight}\ \chi\ Z
 \Longleftrightarrow
 L(\sigma_*+i\log Z,\chi)=0.
\]
The finite declaration \texttt{ThreeDZero} remains the truncated
ordinate-level focal closure used by the finite pencil proofs.

Two exclusion theorems make the image statement exact.  First,
\texttt{not\_scaleBalanced\_of\_ne} says that an abscissa different from
\(\sigma_*\) is not a native carrier state.  Second,
\texttt{helix\_no\_offline\_cancellation} says that an off-circle Cayley
ledger entry cannot be produced by a weld-fixed helix event, while
\texttt{projection\_bijective\_loss\_ledger} says that the height-to-ledger
map is injective on the carrier image.  Projection can therefore neither
move a native cancellation away from the midpoint nor merge two carrier
events into a counterfeit cancellation.  An off-midpoint analytic
vanishing, if presented to this 3D representation, has no scale-balanced,
drift-free carrier source.

This also separates the exact logical role of
\texttt{ThreeD\_crossings\_are\_real\_zeros}.  Its clause
\(\rho=\texttt{carrierPointAtHeight}(Z)\) is the
\emph{coverage/identification}
assertion that a classical analytic zero belongs to the carrier image.  It
is not the derivation of that image's real coordinate: (MP1)--(MP2) and the
exclusion theorems have already done that without quantifying over analytic
zeros.  Once coverage is granted, taking real parts is of course immediate;
the substantive interface question is whether every analytic zero is
represented by a native focal event.  The paper will use
\texttt{carrierPoint} only in this downstream, normal-form sense.  Lean makes
the separation executable by defining the midpoint-free statement
\texttt{NativeCarrierCoverage} and proving
\[
 \texttt{ThreeD\_crossings\_are\_real\_zeros}(\chi)
 \leftrightarrow
 \mathrm{GRH}(\chi)\wedge\texttt{NativeCarrierCoverage}(\chi).
\]
Consequently the legacy equality premise is not cited as evidence for the
coverage direction; an unconditional closure must discharge the midpoint-free
coverage proposition (or the fixed-operator spectral-capture interface)
directly.

\subsection{The chart: \texorpdfstring{\(y=\log Z\) and
\(Z=e^{y}\)}{y = log Z and Z = e\^{}y}}\label{sec:grh-chart}

The readout just constructed does not live at the carrier's scale. The
identity between the scales is one equation read in both directions:
downward, the chart ordinate is the logarithm of the carrier height,
\(y=\log Z\); upward, the carrier height is the exponential of the chart
ordinate, \(Z=e^{y}\). A carrier state at height \(Z>0\) is read at the
represented point \(\tfrac12+iy\) (\texttt{HarmonicPencilCell.reprPoint});
the phasor of the integer \(n\), born at height \(n\), is charted at
ordinate \(\log n\). The compression this induces is the single most
important fact about the relation between the two representations. At the
first zero of \(\zeta\), the chart reads \(y=14.1347\ldots\); the carrier
height is \(Z=e^{y}\approx1.38\times10^{6}\), and the realized bank below
that height already holds about \(1.38\) million integer phasors per lane.
In general the chart ordinate \(y\) summarizes a bank of
\(\lfloor e^{y}\rfloor\) phasors in one complex number. The 1D representation
is not a smaller copy of the 3D representation; it is a logarithmically
compressed summary of the same function, and the scale mismatch between the
two charts is itself a theorem-level object (\S\ref{sec:carrier-S}).

The abscissa tells the same story from the other side.  The carrier's
arclength law first derives the unique scale-balanced exponent
\(\sigma_*=\tfrac12\), and conservation then fixes every native mode at zero
radial drift (\S\ref{sec:grh-midpoint-provenance}).  Riemann's substitution
\(s=\tfrac12+ti\) is the analytic chart of that derived weld axis, and the
represented point \(\sigma_*+iy=\tfrac12+iy\) simply reads height along it.
The ordinate compresses by a logarithm while the abscissa is pinned by a
geometric scale law and its no-drift conservation --- which is why every
subsequent argument in this part is an argument about heights.  The midline
has already been located before the chart is defined.

\subsection{Vanishing is focal alignment}\label{sec:grh-focal}

A zero, on the carrier, is not a value creeping down to nothing; it is an
alignment event. At a marked height the accumulated bank achieves exact
harmonic cancellation: the P and M lanes balance, the \(\mu_6\) cells close in
\(\mathbb Z[\zeta_6]\), and the residual carries no independent constant mode
--- the normalized cell residual factors as \(\det(L^{\mathrm c})/A=
(\lambda-\mu)B\) with \(\lambda-\mu\neq0\), so the residual vanishes exactly
when the signed closure channel does (Theorem~\ref{thm:g6-mode-bridge}).
The event is registered by rank at both exact levels.  At physical helix
height \(Z>0\), with logarithmic ordinate \(y=\log Z\), the canonical finite
bank closes exactly iff its spatially lifted harmonic Gram matrix drops rank
(\texttt{threeDZeroAtHeight\_twoGram3D\_rankDrop}).  This finite event is the
truncated detector \texttt{ThreeDZeroAtHeight}.  The completed fibre event is
\texttt{CompletedThreeDZeroAtHeight}; Lean proves
\[
 \texttt{CompletedThreeDZeroAtHeight}\ \chi\ Z
 \quad\Longleftrightarrow\quad
 L(\sigma_*+i\log Z,\chi)=0
\]
by \texttt{completedThreeDZeroAtHeight\_iff\_L\_zero}.  The theorem
\[
\texttt{completedHarmonicGram3DAtHeight\_rankDrop\_iff\_L\_zero}
\]
proves directly that its completed spatial harmonic Gram lift drops rank exactly at
that represented \(L\)-zero.  Thus the
3D crossing is focal cancellation in the completed \(L\)-function fibre at
physical height \(Z\); it is not inferred from a 1D chart criterion.  The
finite bank theorem is retained as its internal truncation-level detector and
is not substituted for the completed identity.

\medskip
\noindent\textbf{Note (Frobenius determinant one).} The conjugate eigenstate pair the alignment
realizes is not free to sit anywhere; it is pinned by an algebraic constraint. On the carrier the
local Frobenius acts as a \emph{similitude of determinant one}: its chiral block has \(\det=1\)
(\texttt{RequestProject/FrobeniusSimilitude.lean}, \texttt{frobeniusBlock\_det\_one}), so its two
eigenvalues multiply to one --- a reciprocal pair \(\{\alpha,\alpha^{-1}\}\). With the class
unitary they lie on the unit circle as a conjugate pair \(\{e^{i\theta},e^{-i\theta}\}\), and the
eigenphases \(\pm\theta\) are real: the algebraic reason the eigenstate pair is symmetric about the
weld axis rather than drifting off it. This is the carrier's realization of the classical
unit-circle normalization of the Satake--Frobenius class; on the analytic side it is Deligne's
purity, which places the normalized Frobenius eigenvalues on the unit circle \cite{Deligne}, with
the trivial similitude factor \(\det=1\) keeping them a conjugate pair. The one setting where this
mechanism is already a completed theorem is the function-field analog: Weil's Riemann hypothesis for
curves over a finite field \(\mathbb F_q\) is precisely the statement that the Frobenius eigenvalues
on the Jacobian have absolute value \(q^{1/2}\) and pair as \(\{\alpha,q/\alpha\}\), which after
normalization sit as conjugate pairs \(\{\beta,\beta^{-1}\}\) of determinant one on the unit circle
--- forcing the zeros of the congruence zeta function onto \(\Re s=\tfrac12\) \cite{Weil}; Deligne's
proof of the Weil conjectures extends this purity to smooth projective varieties of every dimension
\cite{Deligne}. There the Frobenius \emph{is} the operator whose spectrum pins the zeros to the
critical line, and determinant-one unitarity is the entire content of that conclusion --- the
carrier's \(\det=1\) block is the local, archimedean shadow of the same mechanism. We claim the
analog, not the classical case: the finite-field statement is an unconditional theorem, the
classical one is not, and importing the former as a proof of the latter is exactly the elision this
part refuses. The determinant-one fact is
a structural property of the representation, proven on the carrier; it does not by itself adjudicate
the naming decisions of \S\ref{sec:grh-overview} --- it fixes where the paired eigenphases live, not
which of \((\text{Z-3D})\)/\((\text{Z-1D})\) names the zero.

\subsection{The two Gram operators}\label{sec:grh-operators}

Two operators read the bank, and they divide the labor.

\emph{The Gram harmonic pencil} is the finite detector: the \(2\times2\)
pencil
\(L^{\mathrm c}=\left[\begin{smallmatrix}A&B\\ \mu A&\lambda
B\end{smallmatrix}\right]\)
on the signed channels, whose determinant vanishes --- whose rank drops ---
exactly at the alignment events (\S\ref{sec:grh-focal}). It is the instrument
that marks native focal-event heights on the carrier.

\emph{The local von Neumann fibre} is the coordinate-reading object. At each
physical 3D height \(Z>0\), its spectral parameter is the logarithmic ordinate
\(\gamma=y=\log Z\), and it is multiplication by \(\log Z\) on the
one-dimensional fibre,
\[
\texttt{vonNeumannOp}(\log Z)=(\log Z:\mathbb C)\cdot\mathrm{id},
\]
symmetric --- hence self-adjoint, with real spectrum \(\{\log Z\}\) --- as a
theorem, not an assumption (\texttt{vonNeumannOp\_isSymmetric},
\texttt{RequestProject/ClosedForm.lean}). The spectral parametrization sends
a real eigenvalue \(\mu\) to the represented point
\(\sigma_*+i\mu=\tfrac12+i\mu\), where \(\sigma_*\) has already been fixed
by (MP1).

Because a separate scalar operator is supplied for each \(\gamma\), this family is
not itself the one fixed operator required by the global Hilbert--P\'olya
interface.  Its exact role is to verify the local kernel and the reality of an
already supplied height.

The family is assembled, without changing any local proof, by the fixed
ambient carrier operator \(A^{(3)}_{p,r}\) above.  Its basis is indexed by all
geometric states $X=\gamma_{p,r}(y)$, independently of crossings, and
\(A^{(3)}_{p,r}\delta_X=y\delta_X\).  A completed 3D crossing
$e=(Z_e,h_e)$ selects the state $X_e=\gamma_{p,r}(\log Z_e)$ and Lean proves
that its physical third coordinate is $Z_e$.  Hence a ``3D eigenvector'' in
the ambient resolvent statements below means this explicit carrier basis state, while
\(\texttt{vonNeumannOp}(\log Z_e)\) computes its local
one-dimensional resolvent factor.

The internal 3D chain runs entirely on these two instruments, and every link
is unconditionally proven.  Finite focal cancellation at ordinate
\(y=\log Z\) is detected by the finite harmonic Gram rank drop.  Completed
focal cancellation at physical height \(Z\) is exactly the analytic
\(L\)-zero of that fibre and is detected by the completed harmonic Gram rank
drop.  The spatial statement
\texttt{completedThreeDZeroAtHeight\_twoGram3D\_rankDrop} supplies both the
completed harmonic rank drop and the von Neumann rank drop at the same
ordinate \(\log Z\).  The carrier point
\(\sigma_*+i\log Z=\tfrac12+i\log Z\) lies on the midline because the area law has
already proved \(\sigma_*=\tfrac12\), while von Neumann reality supplies the
real spectral ordinate and hence zero source drift
(\S\ref{sec:grh-midpoint-provenance}).  The one-line Lean proof
\texttt{simp [carrierPoint]} in \texttt{threeDZero\_on\_midline} is now
backed by \texttt{carrierAbscissa\_eq\_half} and is a
normal-form check, not the provenance of the half. Thus focal
cancellation \(=\) 3D zero \(\to\) harmonic Gram rank drop plus von Neumann
Gram rank drop at \(y=\log Z\): that is the detecting and realizing
instrument of this part, and no step consults the chart.

The same-event assertion is now a single compiled assembly theorem, rather
than an inference made in prose.  The theorem
\texttt{completedThreeDZeroAtHeight\_focal\_rankDrop\_eigenvector\_noRadialDrift}
returns, from a completed 3D event at physical height \(Z\), its focal residual
zero, spatial harmonic-pencil rank drop, nonzero ambient eigenvector with
eigenvalue \(\log Z\), and zero radial drift.  Its last component is proved by
\texttt{carrierPointAtHeight\_noRadialDrift}, which uses only that the readout
is a point of the area-law carrier.  Moreover,
\texttt{PrincipalContourNative3DCertificate.analyticFiber\_eq\_carrierFiber}
identifies the literal three-dimensional \(L\)-fiber at \(\rho\) with the
fiber at that certificate's marked carrier height, and the companion theorem
\texttt{.noRadialDrift} transports the geometric no-drift result to the
original analytic parameter.  Thus no-drift is an output of native helix
realization, not an extra condition attached to a zero.
More precisely, the compiled equality
\texttt{eventHeight\_eq\_focalCancellationHeight} identifies the event's
physical height with the literal analytic fiber's stored cancellation height.
The theorem \texttt{harmonicRankDrop\_at\_focalCancellationHeight} rewrites the
completed pencil by this equality before applying its rank-drop certificate;
it does not detect a second event having the same ordinate.  Finally,
\texttt{focalCancellation\_rankDrop\_noRadialDrift\_sameHeight} packages the
readout closure, the determinant zero, and the carrier no-drift law at that one
physical height.

The resolvent trace now has two event-typed versions
(\texttt{RequestProject/TwinResolventTrace.lean}).  For a finite
\(E\subset\mathcal E_\chi\), multiplicities \(m_e\), and
\(p(Z)=\texttt{carrierPointAtHeight}(Z)\), Version~I is
\[
 T^{(3)}_{\chi,E}(w)
 =\sum_{e\in E}m_e\,
   \texttt{vonNeumannResolventTrace}(\log Z_e,w)
 =\sum_{e\in E}\frac{m_e}{w-p(Z_e)}.
\]
The compiled ambient bundle
\texttt{threeDZero\_ambientEigenvector\_resolventTrace} states together that every
index \(e\) is a completed 3D zero, that \(\delta_e\ne0\) is its eigenvector
of \(A^{(3)}_{p,r}\), and that the displayed carrier-point trace formula
holds.  This is precisely the ``3D zeros plus 3D eigenvectors'' version.

Version~II reads the same certified events in the ordinate chart:
\[
 T^{(1\leftrightarrow3)}_{\chi,E}(u)
 =\sum_{e\in E}\frac{m_e}{\log Z_e-u}.
\]
The ambient bundle
\texttt{oneDZero\_threeDZero\_ambientEigenvector\_resolventTrace} additionally
proves, for every \(e\),
\[
 L\!\left(p(Z_e),\chi\right)=0,
 \qquad
 \texttt{CompletedThreeDZeroAtHeight}\ \chi\ Z_e,
 \qquad
 A^{(3)}_\chi\delta_e=(\log Z_e)\delta_e.
\]
The first equation is obtained from the completed focal residual by
\texttt{completedThreeDEigenEvent\_L\_zero}; it is not inserted as a second
unrelated zero assumption.  Finally Lean proves the exact chart relation
\[
 T^{(3)}_{\chi,E}(\texttt{lineC}(u))
 =i\,T^{(1\leftrightarrow3)}_{\chi,E}(u)
 \quad(\texttt{completedTraces\_agree}).
\]
The factor \(i\) is \(d(\texttt{lineC})/du\), and the two pole coordinates
are \(p(Z_e)=\sigma_*+i\log Z_e\) and \(\log Z_e\), respectively.  The older
abstract finite-set results \texttt{trace3D\_capture},
\texttt{trace3D\_analyticAt\_off}, \texttt{trace1D\_capture}, and
\texttt{twinTraces\_agree\_on\_window} remain in place; the new theorems
identify the 3D zero and eigenvector data that their abstract ordinate set did
not itself carry.

\subsection{The chart-side registration audit}\label{sec:grh-audit}

One further object appears in the evidence dossiers and deserves an exact
label. Over a finite window \((0,T]\), the Lean development builds a
diagonal matrix on the window's zero-events --- \(\texttt{hpOperator}\ T\),
one axis per unit of multiplicity, eigenvalue the event ordinate. Its name
is historical; its role is not spectral detection. It is a finite
\emph{repackaging of the already-located analytic zero window}, used for
exactly one thing: auditing the chart against the carrier, count and
multiplicity. The audit facts are proven and exact --- the diagonal is
Hermitian by type (\texttt{hpOperator\_isHermitian}); every index is an
exact vanishing (\texttt{eigenheight\_is\_exact\_vanishing}); the chart
reads zero at a height iff an index sits there
(\texttt{chartZero\_iff\_eigenheight}); the diagonal's dimension is the
chart's own count formula \(1+\theta(T)/\pi+S_{\mathrm{mult}}(T)\)
(\texttt{hpDimension\_eq\_registration}); the eigenspace dimension at an
event equals the vanishing order --- the mode space is a graded local
algebra \(\mathbb C[X]/((X-\mathrm{line}\,\gamma)^{\mathrm{ord}})\)
(\texttt{finrank\_modeSpace}), a double zero a two-dimensional crossing;
and the window traces are honest analytic objects
(\texttt{gradedResolventTrace\_eq\_residue\_sum},
\texttt{windowedDiffResolvent\_tendsto}). But because the window is indexed
by already-located analytic zeros, nothing in this subsection participates
in locating an event.  The local von Neumann fibre of
\S\ref{sec:grh-operators} likewise verifies an event after its height is
supplied.  The fixed global operator interface is stated separately below.

\subsection{The Hilbert--P\'olya setup we use}\label{sec:grh-hp-setup}

The classical suggestion runs: exhibit a self-adjoint operator whose spectrum
gives the ordinates of the nontrivial zeros; self-adjointness forces the
spectrum real; real ordinates place the zeros on the midline. It is worth
saying plainly that this is not an exotic reformulation: Riemann's own
sentence --- \emph{alle Wurzeln reell} --- already asks for a real spectrum,
and real spectrum is precisely what self-adjointness delivers. The
Hilbert--P\'olya suggestion is Riemann's phrasing transposed into operator
language, which is why we treat it as the natural frame for the worked
example rather than as one strategy among many.

The present development has three operator levels, and they must be kept
separate.  The proved scalar fibre
\(\texttt{vonNeumannOp}(\gamma)=\gamma\,\mathrm{id}_{\mathbb C}\) is a local
detector attached after a real height \(\gamma\) is supplied.  Its symmetry
and kernel equation certify the coordinate of that already marked event;
they do not by themselves provide one operator whose spectrum is the event
set.  The fixed \(A^{(3)}_\chi\) supplies the next level: one symmetric
operator on all completed 3D crossing states of \(\chi\), with explicit
eigenvectors and the two resolvent traces above.  Because its index is the
completed-event subtype, it is a valid spectral realization of the 3D event
ontology, not a discovery theorem for analytic zeros not yet coupled to that
ontology.  The final analytic-coverage interface is
\texttt{grh\_of\_selfAdjoint\_resolvent\_capture}: one fixed self-adjoint
element \(a\) must be accompanied by a resolvent-trace identification proving
that every analytic pole parameter lies in \(\operatorname{spectrum}(a)\).
Self-adjointness then supplies spectral reality.  The carrier and chart
theorems below now provide the concrete candidate \(A^{(3)}_\chi\) and its
event traces; applying the final interface to every analytic zero still
requires the coverage/trace identification rather than only the diagonal
assembly.

\section{The theorems}\label{sec:grh-theorems}

The formal home of this part is
\texttt{RequestProject/HelixResolventCapture.lean} together with the
supporting files named below; every statement carries axiom footprint
\(\{\texttt{propext},\texttt{Classical.choice},\texttt{Quot.sound}\}\), no
\texttt{sorry}, no \texttt{axiom}, built in-tree (App.~\ref{app:repro}).
Below, \(\mathrm{GRH}(\chi)\) abbreviates the formal statement ``every
nontrivial zero of \(L(s,\chi)\) lies on the midline''
(\texttt{GRHSpectral.GRH}), stated for every modulus \(N\) and every
Dirichlet character \(\chi\bmod N\). The reading-hypotheses and both GRH
implications are proven at that generality. The trivial character is the
instantiation that lands Mathlib's \texttt{RiemannHypothesis} --- the
mod-\(1\) \(L\)-function \emph{is} \texttt{riemannZeta} --- and a reading
adopted at every modulus yields \texttt{GRHSpectral.GRHComplete}, the
conjecture itself: every Dirichlet \(L\)-function, every modulus, every
character, principal included.

\medskip
\noindent\textbf{Midpoint provenance and exclusion} --- proven before any
zero-set coverage hypothesis:
\begin{itemize}
\item \texttt{windIntegerSite\_radius\_sq\_tendsto} and
\texttt{carrierScaleBalanced\_iff}: arclength forces
\(r_n\asymp\sqrt n\), and \(n^{-\sigma}r_n\) has a positive finite limit
exactly at \(\sigma=\tfrac12\).
\item \texttt{carrierAbscissa\_scaleBalanced},
\texttt{carrierAbscissa\_eq\_half}, and
\texttt{carrierPoint\_eq\_half\_add\_I}: the carrier point is constructed
from the unique scale-balanced abscissa and only afterward rewritten in its
half-unit normal form.  The physical-height readout is separately defined as
\(\texttt{carrierPointAtHeight}(Z)=\texttt{carrierPoint}(\log Z)\), and
\texttt{carrierPointAtHeight\_exp} proves the exact conversion between
physical height and ordinate.
\item \texttt{source\_noDrift}, \texttt{SourceMode.noDrift}, and
\texttt{Helix.no\_radial\_drift\_iff\_half}: norm conservation first gives
\(\Re\lambda=0\) without mentioning a zero or midpoint; after the area-law
normalization, zero radial drift is equivalent to
\(\sigma=\sigma_*=\tfrac12\).
\item \texttt{ThreeDZero = focalClosure} and
\texttt{focalClosure\_iff\_rankDrop}: finite 3D vanishing is defined and
detected without \texttt{carrierPoint}.  At the completed-fibre level,
\texttt{completedThreeDZeroAtHeight\_iff\_L\_zero} identifies focal
cancellation at \(Z\) exactly with
\(L(\texttt{carrierPointAtHeight}(Z),\chi)=0\).
\item \texttt{not\_scaleBalanced\_of\_ne},
\texttt{helix\_no\_offline\_cancellation}, and
\texttt{no\_offCarrier\_representation}: an off-midpoint abscissa is
not a native scale-balanced state; an off-circle ledger cancellation is not
weld-fixed; and no complex point with a different real coordinate is the
image of any real carrier height.  The exact image theorem is
\texttt{exists\_eq\_carrierPoint\_iff}; its positive physical-height form is
\texttt{exists\_pos\_eq\_carrierPointAtHeight\_iff}, with exclusion theorem
\texttt{no\_offCarrier\_physical\_representation}.
\end{itemize}

\medskip
\noindent\textbf{The reduction spine} --- proven, unconditional:
\begin{itemize}
\item \texttt{RH\_of\_GRH\_Trivial\_Char}: \(\mathrm{GRH}(\chi_1)\implies\)
Mathlib's \texttt{RiemannHypothesis}. The hypothesis handles the strip;
Mathlib supplies no-zeros on \(\Re s\ge1\) and the negative-real trivial
zeros via the completed functional equation. The entire target is one
character's GRH.
\item \texttt{focalClosure\_iff\_rankDrop}: exact native finite focal
cancellation is equivalent to harmonic Gram rank drop.  This internal
finite-carrier theorem contains no analytic zero set.
\item \texttt{completedThreeDZeroAtHeight\_iff\_L\_zero}: the completed
3D fibre crosses at physical height \(Z>0\) exactly when \(L\) vanishes at
\(\texttt{carrierPointAtHeight}(Z)=\sigma_*+i\log Z\).  This packages
\texttt{focal\_residual\_zero\_iff\_L\_zero} together with
\texttt{reprPoint\_eq\_carrierPointAtHeight}.  It is an exact pointwise
identity at every represented height; it does not assert that an arbitrary
analytic zero has a representing \(Z\).
\item \texttt{spatialLiftGram\_det},
\texttt{spatialLiftGram\_rankDrop\_iff}, and
\texttt{completedHarmonicGram3DAtHeight\allowbreak\_rankDrop\allowbreak\_iff\allowbreak\_L\allowbreak\_zero}: the
completed harmonic 3D Gram operator rank-drops exactly at the represented
\(L\)-zero.  In addition,
\texttt{completedThreeDZeroAtHeight\_twoGram3D\_rankDrop} states that the harmonic and
von Neumann operators act on the three spatial carrier coordinates as
\(G\otimes I_3\); their determinants are \((\det G)^3\), so the lift
preserves and reflects rank drop and cannot manufacture an extra
cancellation.
\item \texttt{vonNeumannOp\_isSymmetric}: symmetry of the local scalar fibre
operator is a theorem; it supplies coordinate reality after an event has
been coupled to that fibre.
\item \texttt{threeDCrossings\_iff\_grh\_and\_nativeCarrierCoverage}: the
legacy identity premise is exactly the conjunction of the critical-line
conclusion and midpoint-free native coverage.
\item \texttt{grh\_of\_resolvant\_trace\_3D\_real}: for every modulus and
character, \(\textsc{ThreeDRealEvidence}\wedge
\texttt{ThreeD\_crossings\_are\_real\_zeros}(\chi)\implies
\mathrm{GRH}(\chi)\). The evidence component is formally redundant;
\texttt{grh\_of\_threeDCrossings} obtains the conclusion directly from the
equality embedded in the legacy premise.
\item \texttt{grh\_of\_resolvant\_trace\_1D\_correlation}: for every modulus
and character, \(\textsc{OneDCorrelationEvidence}(\chi)\wedge
\texttt{OneD\_zeta\_zero\_correlated}(\chi)\implies\mathrm{GRH}(\chi)\).
The routing proposition is redundant and is itself proved by
\texttt{oneD\_zeta\_zero\_correlated\_proved};
\texttt{grh\_of\_oneDCorrelationEvidence} shows that the chart-to-kernel
premise alone supplies the conclusion.
\item \texttt{grh\_of\_selfAdjoint\_resolvent\_capture}: the global operator
route.  It uses one fixed self-adjoint operator and a trace/capture premise
placing every analytic pole parameter in its spectrum.
\item \texttt{grhComplete\_of\_resolvant\_trace\_3D\_real},
\texttt{grhComplete\_of\_resolvant\_trace\_1D\_correlation}: a reading
adopted at every modulus yields \texttt{GRHComplete} --- the full
conjecture, principal characters included.
\item \texttt{RH\_of\_resolvant\_trace\_3D\_real},
\texttt{RH\_of\_resolvant\_trace\_1D\_correlation}: the trivial-character
instantiations composed with the bridge, landing \texttt{RiemannHypothesis}
directly.
\item \texttt{threeDRealEvidence}, \texttt{oneDChartEvidence}: the two
dossiers inhabited --- the \(16\)-fold and \(18\)-fold conjunctions of
\S\ref{sec:grh-E2} and \S\ref{sec:grh-E1} certified theorem by theorem.
\item \texttt{threeD\_real\_case}, \texttt{oneD\_correlation\_case}: each
reading packaged as one auditable object --- its proven evidence paired with
its displayed conditional premise.  The minimal dependencies are the audit
theorems listed above.
\end{itemize}

\medskip
\noindent\textbf{The two legacy interfaces, rendered formal.} The formal
system speaks the classical frame: its ``nontrivial zero'' is definitionally
the analytic object.  The revised source records the minimal dependency of
each older interface:
\begin{itemize}
\item \texttt{ThreeD\_crossings\_are\_real\_zeros} (\textsc{3D-Zero-Real};
the strong reading with the 3D zero primary): every nontrivial zero \(\rho\)
is represented by a completed fibre crossing at a physical 3D height
\(Z>0\), with
\(\rho=\texttt{carrierPointAtHeight}(Z)=\sigma_*+i\log Z\). The two
spatially lifted Gram pencils then provide their rank drops at ordinate
\(\log Z\). Here \(\tfrac12+i\log Z\) is the proved normal form of the
carrier image, not a fresh
geometric assumption.  Nevertheless, the equality already includes the
critical-line conclusion.  The theorem
\texttt{threeDCrossings\_iff\_grh\_and\_nativeCarrierCoverage} exposes its
exact content; the midpoint-free coverage component is now available as the
separate target \texttt{NativeCarrierCoverage}.
\item \texttt{OneD\_zeta\_zero\_correlated} (\textsc{1D-Zero-Exclusive-Real};
the weak reading): a per-zero disjunction.  Its right branch is only the
analytic zero equation together with \(\gamma=\Im\rho\), so
\texttt{oneD\_zeta\_zero\_correlated\_proved} inhabits it unconditionally.
The actual pointwise coupling is \texttt{OneDCorrelationEvidence}, which
asserts a scalar-fibre kernel at that analytic zero and therefore implies the
critical-line conclusion by itself.
\end{itemize}
These audits do not alter either conditional theorem.  They state which
premise actually supplies the conclusion and prevent the geometric midpoint
theorems, the internal rank-drop correspondence, and the analytic zero-set
identification from being presented as one result.

The coordinate convention in both interfaces is fixed: $Z>0$ is physical
height on the Archimedean helix, $y=\log Z$ is the analytic ordinate, and
$S(t)$ is the global carrier-to-chart coordinate/registration map.  The
pointwise fibre readout and the global census map therefore have distinct
typed roles.

\section{Evidence I: the 3D--1D correlation}\label{sec:grh-E1}

The eighteen facts of \texttt{oneDChartEvidence}, each proven
unconditionally, form a chart-side registration dossier.  They establish
count, jump, residue, multiplicity, and projection identities on the objects
appearing in their statements.  The universal layer is the \(S(t)\) global
coordinate map: it transports the carrier-scale census to the unit-chart
census and registers every analytic event jump.  A global census identity is
not by itself a pointwise bijection between \texttt{ThreeDZero} and arbitrary
Dirichlet \(L\)-zeros.  The separate, character-indexed input
\(\textsc{OneDCorrelationEvidence}(\chi)\) is the operative pointwise
chart-to-kernel leg.
Three of the facts below are \emph{named} engines: \emph{Shannon} --- the collapse loses
radius and angle to the ledger but never counterfeits (\texttt{record\_bijective}); the
\emph{Shannon Cascade} --- an accumulation of stage features is a feature of the limit
(\texttt{limit\_zero\_of\_stage\_accumulation}); and \emph{Riemann's Fold} --- the
projection pipeline, crossing dimensions, lands on the real axis exactly when its retained
coordinate does (\texttt{pipeline\_midpoint\_iff}).

\begin{itemize}
\item \texttt{chartZero\_iff\_eigenheight}
(\texttt{RequestProject/HilbertPolya.lean}). For \(0<\gamma\le T\):
\(\zeta(\tfrac12+i\gamma)=0\iff\gamma\) is an index of the zero-window
diagonal (chart-side audit, \S\ref{sec:grh-audit}). The one-for-one
coincidence, proven --- not asserted --- window by window.
\item \texttt{hpDimension\_eq\_registration}
(\texttt{RequestProject/HilbertPolya.lean}). The window diagonal's dimension
is the chart's own count formula, \(1+\theta(T)/\pi+S_{\mathrm{mult}}(T)\).
The index count and the zero count are the same number because the ledger
says so.
\item \texttt{carrier\_scale\_compensation\_S}
(\texttt{RequestProject/CarrierScaleCompensation.lean}).
\(N_{\pi/3}(e^{t})-N_{1}(e^{t})=S(t)\): the classical \(S(t)\) is the gap
between the native-scale event count and the unit-scale smooth term --- one
procession, two registrations, no mysterious remainder.
\item \texttt{native\_identification$'$}
(\texttt{RequestProject/CarrierScaleCompensation.lean}).
\(N_{\pi/3}(e^{t})\) equals the number of on-line zero ordinates in
\((0,t]\) --- one zero per completed \(\pi/3\) cell.
\item \texttt{S\_jump\_detects\_event}
(\texttt{RequestProject/ResidueJump.lean}). \(S\) has a unit jump at
\(\gamma\) iff \(\zeta(\tfrac12+i\gamma)=0\): detecting a zero and detecting
an \(S\)-jump are the same act.
\item \texttt{count\_hasJump\_iff\_S\_hasJump}
(\texttt{RequestProject/ResidueJump.lean}). The event count and \(S\) have
identical jumps --- their difference is continuous, so \emph{all}
discontinuity lives in \(S\), none in the smooth clock \(\theta/\pi\).
\item \texttt{residue\_eq\_Smult\_jump}
(\texttt{RequestProject/ResidueJump.lean}). At every event,
log-derivative residue \(=\) vanishing order \(=\) ledger jump, with
multiplicity and with \emph{no} simplicity hypothesis.
\item \texttt{record\_bijective} (\texttt{RequestProject/RoundTrip.lean}).
\(\texttt{record}\,f=(\text{height},(\text{radial},\text{phase}))\) is a
bijection with explicit inverse \texttt{reconstruct}: the loss ledger made
literal (\S\ref{sec:grh-projection}) --- \emph{Shannon}: the collapse never counterfeits.
\item \texttt{radial\_lost} (\texttt{RequestProject/RoundTrip.lean}).
Without the ledger the geometric projection is provably non-injective ---
the ledger does real work; the faithfulness results are not vacuous.
\item \texttt{projection\_bijective\_loss\_ledger}
(\texttt{RequestProject/SourceHolonomy.lean}). The chart chain
\(y\mapsto1-(\tfrac12+yi)^{-1}\) is injective: distinct helix events write
distinct ledger entries --- the data-processing inequality made literal.
\item \texttt{pipeline\_midpoint\_iff}
(\texttt{RequestProject/RoundTrip.lean}). The end-to-end
3D\(\to\)height\(\to\)encode\(\to\)1D pipeline reads \(\tfrac12\) exactly
when the retained 3D coordinate is \(\tfrac12\): the midline address
survives projection intact, and only the midpoint maps there --- \emph{Riemann's Fold}.
\item \texttt{no\_radial\_drift\_on\_helix}
(\texttt{RequestProject/SourceHolonomy.lean}). The Euler-factor phasor's
radius depends only on \(\Re s=\tfrac12\), not on height: the
height\(\to\)circle map is a faithful re-coordinatization, all
distinguishing information in the phase the ledger keeps.
\item \texttt{midpoint\_entry\_on\_circle}
(\texttt{RequestProject/SourceHolonomy.lean}). Every helix height charts to
a unit-circle ledger entry (\(\|1-(\tfrac12+yi)^{-1}\|=1\)): line-events are
recorded as circle-points.
\item \texttt{faithful\_projection\_zeta}
(\texttt{RequestProject/Faithfulness.lean}). The projected fiber is real on
the midline and vanishes exactly at the genuine on-line zeros: the collapse
neither hides a zero nor manufactures one.
\item \texttt{euler\_maclaurin\_dirichlet}
(\texttt{RequestProject/EulerMaclaurinDirichlet.lean}). In the strip, the
finite Dirichlet sum with boundary correction approximates \(\zeta\) with
error uniform in the truncation --- the quantitative backbone of the chart
readout.
\item \texttt{transfer\_tendsto}
(\texttt{RequestProject/TransferContinuation.lean}). A polynomial bound on
the primitive continues the series past the trivial abscissa: the readout
limit exists from a growth estimate alone, with no appeal to zero locations.
\item \texttt{limit\_zero\_of\_stage\_accumulation}
(\texttt{RequestProject/ExportAdapter.lean}). A zero accumulating across the
finite stages is a zero of the limit \(\xi\)-section: the collapse does not
counterfeit features --- the \emph{Shannon Cascade} converse.
\item \texttt{windowedDiffResolvent\_tendsto}
(\texttt{RequestProject/DifferencedResolvent.lean}). The windowed
differenced-resolvent trace converges as \(T\to\infty\) to an absolutely
convergent limit: the chart-side resolvent is an honest analytic object,
unconditionally.
\end{itemize}

\section{Evidence II: internal three-dimensional carrier facts}
\label{sec:grh-E2}

The sixteen facts of \texttt{threeDRealEvidence} are each proven
unconditionally, but their conjunction is an internal-support dossier, not a
theorem identifying every analytic zero with a \texttt{ThreeDZero}.  The
facts concern local scalar-fibre symmetry, finite carrier cancellation,
completed critical-line readouts, chart-indexed registration objects,
arithmetic winding, and general limit principles.  Their scopes are recorded
individually below.  Three named engines organize the dossier:
\emph{Feynman's Quiver} --- the crossings are prime-generated by unique factorization, the
phasor bank is summable, the heights are distinct, and no focus splits; \emph{Deligne's
Pairs} --- the self-reciprocal det-one local pairing, purity at the finite places; and
\emph{Shannon Projection Dominance} --- the source dominates every projection, so the limit
holds no feature the stages lack.

\begin{itemize}
\item \texttt{vonNeumannOp\_isSymmetric}
(\texttt{RequestProject/ClosedForm.lean}). The fibre operator
\(\gamma\cdot\mathrm{id}\) is symmetric, hence self-adjoint with real
spectrum \(\{\gamma\}\).  This is a local coordinate detector indexed by an
already supplied real \(\gamma\), not the fixed global operator of
\texttt{grh\_of\_selfAdjoint\_resolvent\_capture}.
\item \texttt{carrierThreeDOperator\_isSymmetric},
\texttt{carrierThreeDOperator\_eigenvector}, and
\texttt{carrierThreeDOperator\_completedEvent\_eigenvector}
(\texttt{RequestProject/ThreeDFocalEvent.lean}).  For fixed nondegenerate geometric carrier
parameters $(p,r)$ with $p\ne0$, one real-diagonal operator acts on all finitely supported
carrier states $X=\gamma_{p,r}(y)$; its type mentions no event predicate.
Every geometric state supplies a nonzero basis eigenvector with eigenvalue
$y$, and a completed crossing at $Z_e$ is proved to select the state with
$y=\log Z_e$ and third coordinate $Z_e$.
\item \texttt{threeDZero\_ambientEigenvector\_resolventTrace} and
\texttt{oneDZero\_threeDZero\_ambientEigenvector\_resolventTrace}
(\texttt{RequestProject/TwinResolventTrace.lean}).  The first bundle contains
3D crossing certificates, their fixed-operator eigenvectors, and the 3D pole
formula.  The second adds the analytic \(L\)-zero readout for each same event,
the ordinate trace, and the exact carrier--ordinate chart relation.
\item \texttt{eigenheight\_is\_exact\_vanishing}
(\texttt{RequestProject/HilbertPolya.lean}). Every index of the chart-side
window satisfies \(\zeta(\tfrac12+i\gamma)=0\) exactly: the audit's indices
are genuine events, not approximations (\S\ref{sec:grh-audit}).
\item \texttt{hpOperator\_isHermitian}
(\texttt{RequestProject/HilbertPolya.lean}). The zero-window registration
diagonal is Hermitian --- reality by type (chart-side audit,
\S\ref{sec:grh-audit}).
\item \texttt{finrank\_modeSpace}
(\texttt{RequestProject/GradedModeDictionary.lean}). The mode space at an
event is a graded local algebra of dimension exactly the vanishing order: a
double zero is a two-dimensional crossing; spectral multiplicity \(=\)
analytic order.
\item \texttt{gradedResolventTrace\_eq\_residue\_sum}
(\texttt{RequestProject/GradedModeDictionary.lean}). The 3D-mode resolvent
trace equals the explicit-formula sum of log-derivative residues: the trace
formulas coincide, not just the locations.
\item \texttt{focal\_residual\_zero\_iff\_L\_zero}
(\texttt{RequestProject/FocalResidualVanishes.lean}).
\(D_{\mathrm{cell}}=V_\chi\Phi_\chi\) with \(V_\chi\neq0\) and
\(\Phi_\chi=(\pi/3)\,L\): the cell channel vanishes iff the \(L\)-value
does, residue-free, at exact \(\pi/3\) amplitude at the represented
critical-line point.  This completed-cell statement is distinct from the
finite-bank definition of \texttt{ThreeDZero}.
\item \texttt{fiber\_accumulates\_to\_L}
(\texttt{RequestProject/Faithfulness.lean}). For \(\chi\neq1\) and
\(\Re s>0\) the partial sums converge to \(L(s,\chi)\) across the whole
strip.  This is value convergence at each stated \(s\); transferring zero
sets requires the separate uniform-convergence/isolation hypotheses used by
the limit theorems.
\item \texttt{threeD\_exhaustive}
(\texttt{RequestProject/SourceHolonomy.lean}). Every zero event on the real
height ray is sourced --- automatic because reals are conjugation-fixed.  Its
domain is the real-height event space; it does not prove that every complex
analytic zero lies in that domain.
\item \texttt{threeD\_metric\_no\_zeros}
(\texttt{RequestProject/SourceHolonomy.lean}). The 3D cup form is definite:
no nonzero state is null, so vanishing is a property of \emph{readings},
never of states --- a zero cannot hide as a degenerate state.
\item \texttt{windFromPrimes\_mul}
(\texttt{RequestProject/Origination.lean}). The winding is multiplicative
over factorization, with no \(\Re\!\log\): the fundamental theorem of
arithmetic realized as a homomorphism into the circle --- carrier and
\(L\)-function are two prime-built readings of one arithmetic (\emph{Feynman's Quiver},
generation).
\item \texttt{heights\_distinct}
(\texttt{RequestProject/NoDoubleCancellation.lean}). Distinct sites carry
distinct eigen-heights: unique factorization as a spectral
non-degeneracy (\emph{Feynman's Quiver}, uniqueness).
\item \texttt{cancellation\_modes\_subsingleton}
(\texttt{RequestProject/NoDoubleCancellation.lean}). At any spectral
parameter, at most one basis mode cancels: no split focus --- a crossing
marks one ordinate, not a superposition (\emph{Feynman's Quiver}, no split focus).
\item \texttt{localPoly\_reciprocal}
(\texttt{RequestProject/FiniteWeightFiber.lean}). The local numerator is
self-reciprocal: the per-place functional equation whose global assembly has
fixed axis \(\Re s=\tfrac12\) --- the midline the crossings sit on (\emph{Deligne's Pairs}).
\item \texttt{fiber\_det\_one}
(\texttt{RequestProject/FiniteWeightFiber.lean}). \(\prod_i\lambda_i=1\)
under the duality involution: the modulus ledger the reciprocal law
consumes (\emph{Deligne's Pairs}).
\item \texttt{frobeniusBlock\_det\_one}
(\texttt{RequestProject/FrobeniusSimilitude.lean}). The chiral block
\(\mathrm{diag}(z,\bar z)\) has determinant \(1\): the Frobenius similitude
is a pure rotation --- Satake weights stay on the unit circle, eigenphases
stay real (\emph{Deligne's Pairs}).
\item \texttt{limit\_dominance}
(\texttt{RequestProject/LimitDominance.lean}). Hurwitz-type dominance: the
under local uniform convergence and isolation, a limit zero has nearby stage
zeros eventually.  This gives proximity, not equality with a fixed finite
stage or a proof that the nearby zeros lie on the carrier image.
\end{itemize}

\section{Proof A: GRH under (Z-3D)}\label{sec:grh-proofA}

\begin{theorem}[The identity route]\label{thm:grh-routeA}
Unconditionally:
\begin{enumerate}[label=\textup{(\roman*)},leftmargin=*,nosep]
\item \(\textsc{ThreeDRealEvidence}\) holds
(\texttt{threeDRealEvidence});
\item for every modulus \(N\) and character \(\chi\bmod N\),
\[
\begin{gathered}
\textsc{ThreeDRealEvidence}\ \wedge\\
\texttt{ThreeD\_crossings\_are\_real\_zeros}(\chi)
\ \Longrightarrow\ \mathrm{GRH}(\chi).
\end{gathered}
\]
Lean: \texttt{grh\_of\_resolvant\_trace\_3D\_real};
\item the reading adopted at every modulus yields the full conjecture
\texttt{GRHComplete}.\\
Lean: \texttt{grhComplete\_of\_resolvant\_trace\_3D\_real};
\item at the trivial character the premises imply
\texttt{RiemannHypothesis}.\\
Lean: \texttt{RH\_of\_resolvant\_trace\_3D\_real}.
\end{enumerate}
\end{theorem}

\begin{proof}[Proof A]
Adopt (Z-3D): ``nontrivial zero'' names the crossing --- the focal-alignment
event of \S\ref{sec:grh-focal}.  Two logically distinct facts are then
composed.  First, before this adoption, the arclength area law and
conservative transport have proved that the native, scale-balanced carrier
image is exactly
\(\{\sigma_*+i\log Z:Z>0\}\), equivalently
\(\{\sigma_*+iy:y\in\mathbb R\}\), with
\(\sigma_*=\tfrac12\), and that its ledger cannot encode an off-image
cancellation (\S\ref{sec:grh-midpoint-provenance}).  Second, (Z-3D) supplies
the coverage/identity clause that each classical nontrivial zero is one of
those native focal events.  This second clause is the explicitly displayed
hypothesis in \textup{(ii)}; it is not used in the derivation of
\(\sigma_*\).

For each covered event at physical height \(Z\), the two spatial Gram
theorems provide the harmonic and von Neumann rank drops at the same
ordinate \(y=\log Z\).  The fibre operator is
self-adjoint by theorem
(\texttt{vonNeumannOp\_isSymmetric}), and the represented coordinate is the
already derived \(\sigma_*+i\log Z\).  Hence taking real parts gives
\(\Re\rho=\sigma_*=\tfrac12\).  With the evidence of
\S\ref{sec:grh-E2} --- exact vanishing at the eigenheights, multiplicity
equal to order, no split focus, definite metric, exhaustion of the real ray
--- theorem \textup{(ii)} closes \(\mathrm{GRH}(\chi)\) under the stated
coverage hypothesis; \textup{(iii)} and \textup{(iv)} are its all-modulus
and trivial-character specializations.

Equivalently, the newly assembled operator records the same step without a
family-only ambiguity: the covered crossing is an index
\(e\in\mathcal E_\chi\), its explicit state \(\delta_e\) is an eigenvector of
\(A^{(3)}_\chi\) with real eigenvalue \(\log Z\), and Version~I of the
resolvent trace has its pole at the carrier point of that same event.  This
augments the retained proof; it does not replace its displayed coverage
hypothesis.
\end{proof}

\medskip
\noindent\textbf{Formal dependency audit for Proof A.}  The proof text above
is retained.  The compiler-level implication is sharper: the first evidence
component is unused, because the equality
\(\rho=\texttt{carrierPointAtHeight}(Z)\) inside the second component already yields
the real-part conclusion.  The theorem
\texttt{threeDCrossings\_iff\_grh\_and\_nativeCarrierCoverage} records this
exactly.  Thus Proof~A is a valid conditional deduction, while its legacy
identity premise must not be advertised as an independently proved
analytic-zero coverage theorem.  The geometric result established before
that premise is the location of the native carrier image.

\section{Proof B: GRH under (HP-W)}\label{sec:grh-proofB}

\begin{theorem}[The correlation route]\label{thm:grh-routeB}
Unconditionally:
\begin{enumerate}[label=\textup{(\roman*)},leftmargin=*,nosep]
\item \(\textsc{OneDChartEvidence}\) holds
(\texttt{oneDChartEvidence});
\item for every modulus \(N\) and character \(\chi\bmod N\),
\[
\begin{gathered}
\textsc{OneDCorrelationEvidence}(\chi)\ \wedge\\
\texttt{OneD\_zeta\_zero\_correlated}(\chi)
\ \Longrightarrow\ \mathrm{GRH}(\chi).
\end{gathered}
\]
Lean: \texttt{grh\_of\_resolvant\_trace\_1D\_correlation};
\item the reading adopted at every modulus yields the full conjecture
\texttt{GRHComplete}.\\
Lean: \texttt{grhComplete\_of\_resolvant\_trace\_1D\_correlation};
\item at the trivial character the premises imply
\texttt{RiemannHypothesis}.\\
Lean: \texttt{RH\_of\_resolvant\_trace\_1D\_correlation}.
\end{enumerate}
\end{theorem}

\begin{proof}[Proof B]
Adopt (HP-W): the criterion demands that the operator's eigenvalues
\emph{coincide with} the zeros. The hypothesis in \textup{(ii)} routes each
zero through one of two forms: a crossing at a positive physical 3D height,
or a 1D chart
report at its ordinate.  In the first form the crossing is a completed focal
event at physical height \(Z>0\), read at ordinate \(y=\log Z\). In the latter
form, \(\textsc{OneDCorrelationEvidence}(\chi)\)
is a proof input: it turns that chart report into the nonzero von Neumann
kernel at the reported point. Section \ref{sec:grh-E1} supplies the compiled
chart-side dossier --- the window bijection, registration identity,
jump-for-residue agreement with multiplicity, and projection ledger. Both
branches then close by the same von Neumann reality (the case split inside
\texttt{grh\_of\_resolvant\_trace\_1D\_correlation}), whence
\(\mathrm{GRH}(\chi)\) at every modulus, \textup{(iii)} assembles the full
conjecture, and \textup{(iv)} lands the classical statement. The weak
reading consumes nothing stronger than the strong one --- the same
self-adjoint mechanism closes both branches.

For a chart report that is coupled to a completed crossing, Version~II makes
all three records simultaneous in Lean: the 1D equation
\(L(p(Z_e),\chi)=0\), the completed 3D event certificate, and the 3D
eigenvector \(\delta_e\) of \(A^{(3)}_\chi\).  The identity
\(T^{(3)}(\texttt{lineC}(u))=iT^{(1\leftrightarrow3)}(u)\) is the compiled
correlation between those records.  The pointwise premise in the legacy proof
still performs the routing of an arbitrary chart report into this event
space.
\end{proof}

\medskip
\noindent\textbf{Formal dependency audit for Proof B.}  The proof text above
is retained.  The source now proves
\texttt{oneD\_zeta\_zero\_correlated\_proved}: the disjunctive routing
proposition is automatic because its right branch records only
\(L(\rho,\chi)=0\) and chooses \(\gamma=\Im\rho\).  Conversely,
\texttt{grh\_of\_oneDCorrelationEvidence} proves that the operative evidence
alone gives the critical-line conclusion, since it asserts a scalar-fibre
kernel at each analytic zero.  The global \(S(t)\) coordinate law explains
the count, jump, and multiplicity registration used by this route; the
pointwise chart-zero-to-kernel assertion remains its separate premise.

\section{Riemann Hypothesis corollaries at \texorpdfstring{\(\chi=1\)}{chi=1}}
\label{sec:grh-rh}

The mod-\(1\) \(L\)-function is \texttt{riemannZeta}, so at the trivial
character each conditional implication lands the classical statement with
the displayed premise unchanged:

\begin{corollary}[RH, identity reading]\label{cor:grh-rhA}
\(\textsc{ThreeDRealEvidence}\wedge
\texttt{ThreeD\_crossings\_are\_real\_zeros}(\chi_1)\)
implies \texttt{RiemannHypothesis}
(\texttt{RH\_of\_resolvant\_trace\_3D\_real}, proven).
By the audit above, the identity premise itself contains the
critical-line conclusion as well as native coverage.
\end{corollary}

\begin{corollary}[RH, correlation reading]\label{cor:grh-rhB}
\(\textsc{OneDCorrelationEvidence}(\chi_1)\wedge
\texttt{OneD\_zeta\_zero\_correlated}(\chi_1)\)
implies \texttt{RiemannHypothesis}
(\texttt{RH\_of\_resolvant\_trace\_1D\_correlation}, proven).
The operative \(\textsc{OneDCorrelationEvidence}(\chi_1)\) bridge alone is
the pointwise premise that supplies the critical-line conclusion.
\end{corollary}

Both factor through the same bridge, \texttt{RH\_of\_GRH\_Trivial\_Char}:
one character's GRH is RH.

\section{Conclusion}\label{sec:grh-conclusion}

The corrected dependency structure has three levels.  First, Archimedean
helix geometry derives the unique scale-balanced, drift-free carrier
abscissa; Lean constructs \texttt{carrierPoint} from that abscissa, proves
its half-unit normal form afterward, and reads physical height through
\(\texttt{carrierPointAtHeight}(Z)=\texttt{carrierPoint}(\log Z)\).  Second,
the finite harmonic pencil detects truncated focal cancellation exactly, the
completed harmonic pencil detects the completed \(L\)-fibre cancellation,
and the local scalar fibre verifies its ordinate \(y=\log Z\).  The fixed
operator \(A^{(3)}_\chi\) then assembles all completed crossings into one
symmetric event-state space with explicit basis eigenvectors.  Third,
identifying the analytic zero set with those native events requires a
pointwise coverage or spectral-capture theorem.

The global \(S(t)\) coordinate map supplies exact census, jump, residue, and
multiplicity registration between carrier and chart scales.  The proved
supporting facts include the two explicit completed-event resolvent versions:
3D zeros with 3D eigenvectors, and 1D zero readouts with the same 3D zeros and
3D eigenvectors.  The legacy identity
conditional and correlation conditional also remain valid deductions, with
their proofs retained above; the new Lean audit states exactly which premise
contains the critical-line conclusion in each route.  An analytic
Hilbert--P\'olya completion can now target the fixed
\(A^{(3)}_\chi\): it must prove the analytic resolvent/spectral capture
required by \texttt{grh\_of\_selfAdjoint\_resolvent\_capture}, or
instantiate that interface with another fixed operator.  The scalar family
\(\gamma\,\mathrm{id}_{\mathbb C}\) is correctly used as the local
diagonal-mode calculator, while \(A^{(3)}_\chi\) is the assembled 3D
event-space operator.
